% Modernised reading — fonds Alexandre Grothendieck, shelfmark 5, pages 1-52.
%
% This is an INTERPRETATION, not a transcription. It re-reads the pages in
% today's notation and vocabulary, states the hypotheses the manuscript leaves
% implicit, and drops the critical apparatus so that the mathematics reads
% straight through. Everything the brackets and underlines carried — a doubtful
% reading, a notation he used differently, a missing passage — moves into a
% footnote.
%
% It is held to being mathematically correct as it stands, not merely faithful
% to the page. Where the manuscript is loose or mistaken, the true statement is
% what appears here, and the footnote says what the page has. In any dispute of
% substance the transcription governs: whoever cites Grothendieck must cite the
% batch-NN.fr.tex files.
%
% Scope: ONE document for the whole shelfmark, its three batches of
% transcription (pages 1-52, all transcribed) read as a single argument. The
% folder is a workshop for the apparatus of crystalline cohomology in 1965-66:
% divided powers (three passes: pages 21-26 and 31-34 abstractly, 38-40 in a
% Z_(p)-algebra, 44-46 on a principal ideal, with the artinian computation of
% pages 4-12 as their concrete test), the sites (page 2 and pages 27-30), the
% Frobenius topos (pages 15-19), and the letter to Deligne of 10.12.1965
% (pages 49-52) that says what all this is for. The inventory's « correspondance
% Berthelot » has no leaf among the 52 pages.
%
% NOTATION fixed once for the folder.
%   Divided powers: gamma_n(x) everywhere. The pages write x^{(n)} and pi^n(x)
%     (21-26), Gamma^n(x) (31-34), pi_m(x) (38-40), pi^{(n)} (44-46),
%     gamma^{p^r} (4-8).
%   Sites: his names DRIFT between page 2 and page 28 — his « cris » of page 2
%     is his « DR » of page 28, and his « Cris » of page 28 is his « inf » of
%     page 2. This reading writes Cris(X/S) (with divided powers, today's
%     crystalline site), Cris^st (its sub-site of objects locally retracting),
%     Inf(X/S) (no divided powers) and Inf^st (his Strat). Table in the
%     conventions section.
%   Coefficients: q_i = (p^{i+1})!/((p^i)!)^p (pages 4-8 and 32); kappa_n
%     (pages 31-32, his c_n); e_n and d_m (page 38, his c_n and d_m); kappa'_n
%     (page 40, his c_n). His one letter c_n is three different numbers.
%   PD envelope: D_A(I) (his Gamma*(A, I), Gamma(A, J)).
%
% Substantive departures from the pages, each footnoted where it occurs:
%   - page 5: « si i = nu » for i = 0 in the induction;
%   - page 12: chif(n) <= (p-1) nu should be (p-1)(nu+1) (harmless);
%   - page 16: Homtop(X, Q) for Homtop(X', Q);
%   - page 18: omega_X^* for omega_Y^* in the base-change map;
%   - page 26: log series with n >= 0 and sign (-1)^n, for n >= 1, (-1)^{n-1};
%   - page 31: v_p(n) for v_p(n!); Gamma^n(xy) = x^n Gamma^n(y) with x in I,
%     y in A, roles swapped;
%   - page 32 (margin): a p-th POWER carrying the coefficient of a COMPOSITE;
%   - page 34: boxed formula not as it should read from 1.5; the true one given;
%   - page 40: the relation p! pi(x) = x^p, necessary, is missing from the
%     list; his question « c_n = 1 mod p ? » is answered here — no,
%     kappa'_n = 1 / prod a_i! mod p (checked by machine for p <= 7, n < 200);
%   - page 44, Example 1: « pi^n in n! A » is NOT sufficient for divided powers
%     on the principal ideal pi A; the condition is pi^n in n! pi A
%     (counterexample: A = Z_(p)<T>, pi = T);
%   - page 2: the « classe de Chern » is read as an obstruction class, and the
%     smooth case needs a QUASI-NILPOTENT integrable connection in char. p.
% The lemmas of pages 5 and 7 (normal form in W^{(nu)} and W^{(nu,r)}) were
% checked by machine for (p, nu) in {(2,1),(2,2),(2,3),(3,1),(3,2),(5,1)} and
% every r: the ideal is exactly the diagonal lattice the page states.
%
% Announced and not established on the pages: that W^{(nu)} carries the
% divided powers (page 4, asserted); that the presentation of page 31 is the
% PD envelope (« à vérifier »); the rigidity H^q(Cris^st/(U,U'), O) = H^q(U', O)
% (« à vérifier »); the vanishing of R^i in the Frobenius topos (page 18-19,
% breaks off); the extension of pi_p to a full divided-power structure
% (page 40, open question). The polynomial-maps note closing on page 42 has no
% beginning among the transcribed pages.
%
% Pass: Opus 5.5 (claude-opus-5-5), 2026-10-03 — first pass, unchecked against the pages by a human.
\documentclass[11pt,a4paper]{article}
\input{../preamble/grothendieck}

\folder{5}
\pages{1}{52}
\dating{1965-1969}
\watermark{Édition de démonstration\\interprétation personnelle de l'œuvre}
\foldertitle{Cristaux et cohomologie de De Rham : généralités, correspondances Deligne et Berthelot (1965-1969) : notes manuscrites (s.d.), lettre (1965). — lecture modernisée du dossier entier}

\begin{document}

\section*{Résumé}

\begin{resume}
Ce dossier est un atelier. On y voit Grothendieck, entre 1965 et 1966,
fabriquer les outils d'une théorie qui n'existe pas encore : une cohomologie
des variétés algébriques en caractéristique $p$ qui soit aussi bonne que la
cohomologie de De Rham l'est sur les nombres complexes. Il l'appellera
\emph{cohomologie cristalline}. La lettre à Deligne du 10 décembre 1965, qui
clôt le dossier, dit en une phrase pourquoi il la faut : « on ne pourra dire
des choses vraiment intéressantes et nouvelles qu'en faisant appel à la
“vraie” cohomologie $p$-adique ».

Le problème est simple à formuler. Sur les nombres complexes, la cohomologie
de De Rham se calcule avec des formes différentielles, et tout repose sur le
fait qu'une forme fermée est localement exacte — le lemme de Poincaré, qui
n'est au fond que l'intégration des séries entières : $x^{n}$ a pour
primitive $x^{n+1}/(n+1)$. En caractéristique $p$, on ne peut plus diviser
par $p$ : $x^{p-1}\,dx$ est fermée et n'a pas de primitive, et la théorie
s'effondre. Pire, les développements de Taylor, qui servent à comparer une
fonction en deux points voisins, font intervenir $x^{n}/n!$, qui n'a plus de
sens.

L'idée qui traverse le dossier est de \emph{garder les quotients
$x^{n}/n!$ comme symboles}, sans plus chercher à diviser. On décrète qu'un
élément $x$ possède des « puissances divisées » $\gamma_n(x)$, qui se
comportent comme $x^{n}/n!$ — elles se multiplient, s'additionnent,
se composent selon les mêmes règles — mais qui ne sont pas obtenues en
divisant. C'est comme si l'on gardait dans sa poche les fractions sans
jamais avoir le droit de les écrire comme des nombres. Une bonne moitié des
feuillets met cette notion au point, et l'éprouve : dans quels anneaux ces
symboles existent-ils, combien de choix y a-t-il, que se passe-t-il quand
on les impose à $p$ lui-même ? La réponse la plus nette tient en une
inégalité, encadrée sur la page 46 : dans un anneau de valuation discrète
d'indice de ramification $e$, l'idéal maximal porte des puissances divisées
si et seulement si $e \leqslant p-1$.

Vient ensuite le cadre géométrique. Plutôt que de travailler sur la variété
$X$ elle-même, on considère tous ses « épaississements » : des objets un peu
plus gros que des ouverts de $X$, dans lesquels ceux-ci sont plongés comme
un point est plongé dans un germe de courbe, et dont l'idéal porte des
puissances divisées. L'ensemble de ces épaississements forme un \emph{site},
au sens de la théorie des topos, et sa cohomologie est la cohomologie
cherchée. Les pages 2 et 27 à 30 dressent la carte des variantes possibles
— avec ou sans puissances divisées, avec ou sans rétraction locale — et
disent lesquelles coïncident quand $X$ est lisse ou en caractéristique zéro.

Troisième pièce : Frobenius, l'élévation à la puissance $p$, qui est en
caractéristique $p$ un endomorphisme de tout. Les pages 15 à 19 construisent,
pour deux morphismes de topos parallèles, un nouveau topos dont les objets
sont des faisceaux munis d'une flèche « de Frobenius » — c'est le cadre où
vivront ce qu'on appelle aujourd'hui les $F$-cristaux.

La lettre à Deligne, enfin, regarde plus loin. Elle propose que, sur un corps
$p$-adique, la cohomologie étale et la cohomologie de De Rham soient liées
comme elles le sont sur $\mathbf{C}$ — et aussitôt elle doute : « j'ai de
grands doutes que les choses marchent aussi simplement que ça ». Ce doute était
juste. Ce qu'il pressent là deviendra, quinze ans plus tard, la théorie de
Hodge $p$-adique de Fontaine, où l'on découvre qu'il faut un anneau de
« périodes » pour comparer les deux cohomologies.

Il faut lire ce dossier comme on regarde des esquisses : plusieurs feuillets
reprennent la même construction sous trois notations différentes, deux
calculs s'interrompent sur une question, une conjecture s'énonce et se
rétracte dans la même page. C'est le moment où les objets n'ont pas encore
de nom — « cris » désigne une chose à la page 2 et une autre à la page 28 —
et c'est pour cela qu'il est instructif.

\keywords{divided powers, PD envelope, divided power algebra, Legendre's formula, Witt vectors, Cohen structure theorem, absolute ramification index, crystalline site, infinitesimal site, stratification, crystal, integrable connection, Cech-Alexander complex, Frobenius, inserter, lax colimit of topoi, base change morphism, F-crystal, polynomial law, Poincaré lemma, quasi-homogeneous singularity, completed module of differentials, Poincaré-Verdier duality, rigid analytic space, p-adic Hodge theory, Cartier isomorphism, conjugate spectral sequence}
\end{resume}

\pagerange{2}{2}
\section*{Le fil du dossier, et les conventions valables pour tout le dossier}

\subsection*{Les stations}

Les cinquante-deux feuillets ne forment pas un texte suivi : ce sont des
liasses, souvent écrites au dos de tapuscrits étrangers, que l'archive a
mises dans l'ordre où elles sont. Dans l'ordre des pages :

\begin{itemize}
\item \textbf{Page 2} — la carte des topos attachés à $X$ (cristallin,
infinitésimal, leurs variantes « stratifiées »), et une classe d'obstruction
pour les faisceaux inversibles.
\item \textbf{Pages 4 à 12} — un calcul : l'anneau artinien universel dont
l'idéal maximal porte des puissances divisées nilpotentes d'ordre donné,
avec sa forme normale (pages 4 à 8), la formule de Legendre pour $v_p(n!)$
(page 10), et la convergence de $\pi^{n}/n!$ (page 12).
\item \textbf{Pages 15 à 19} — « Général nonsens sur Frobenius » : le topos
$\mathcal{Q}(\alpha, \beta)$ des faisceaux munis d'une flèche
$\alpha^{*}E \to \beta^{*}E$.
\item \textbf{Pages 21 à 26} — les §§ 1 à 3 d'une rédaction : axiomes des
puissances divisées, algèbre $\Gamma(M)$, enveloppe à puissances divisées
d'un idéal, puissances divisées nilpotentes, $\log$ et $\exp$.
\item \textbf{Pages 27 à 30} — le § 4 : le site de De Rham, sa variante
stratifiée, le carré des topos, et la suite spectrale qui calcule la
cohomologie par les voisinages à puissances divisées de la diagonale.
\item \textbf{Pages 31 à 40} — retour sur l'enveloppe, par générateurs et
relations (pages 31 à 34) ; une liste d'anneaux où vaut le lemme de Poincaré
(page 36) ; le cas d'une $\mathbf{Z}_{(p)}$-algèbre, où tout est déterminé
par $\gamma_p$ (pages 38 et 40).
\item \textbf{Page 42} — la fin d'une note sur les applications polynomiales
homogènes.
\item \textbf{Pages 44 et 46} — « Cristaux » : quand un idéal principal
$\pi A$ porte-t-il des puissances divisées ? Trois exemples, dont
l'inégalité $e \leqslant p - 1$.
\item \textbf{Pages 45 et 47} — au verso des précédentes, deux pages
annotées d'un tapuscrit d'EGA IV, § 18.13.
\item \textbf{Pages 49 à 52} — la lettre à Deligne du 10.12.1965.
\end{itemize}

Logiquement, l'ordre est autre. La théorie des puissances divisées est
exposée trois fois, à trois niveaux : abstraitement aux pages 21 à 26 et 31
à 34 ; pour une $\mathbf{Z}_{(p)}$-algèbre aux pages 38 et 40 ; pour un
idéal principal aux pages 44 et 46. Le calcul des pages 4 à 12 est l'essai
de ces notions sur l'exemple le plus serré qui soit, l'anneau $W$ des
vecteurs de Witt et l'élément $p$. Les pages 2 et 27 à 30 en font le cadre
géométrique, la page 15 y ajoute Frobenius, et la lettre dit à quoi tout
cela doit servir. Ces reprises ne sont pas des répétitions : la page 34 et
la page 44 posent la même question — quelles relations imposer aux
$\gamma_n(\lambda)$ d'un générateur — l'une pour \emph{construire}
l'enveloppe, l'autre pour savoir si l'idéal $\pi A$ \emph{est déjà} à
puissances divisées.

\subsection*{Notations}

\emph{Puissances divisées.} On écrit partout $\gamma_n(x)$, et
$\gamma_0(x) = 1$, $\gamma_1(x) = x$. La page l'écrit de cinq façons :
$x^{(n)}$ et $\pi^{n}(x)$ aux pages 21 à 26, $\Gamma^{n}(x)$ aux pages 31
à 34, $\pi_m(x)$ aux pages 38 et 40, $\pi^{(n)}$ aux pages 44 et 46,
$\gamma^{p^{r}}$ aux pages 4 à 8. Un couple $(A, I, \gamma)$ — un anneau, un
idéal, une structure de puissances divisées sur l'idéal — s'appelle un
\emph{anneau à puissances divisées}.

\emph{Coefficients.} Le dossier emploie une seule lettre, $c_n$, pour trois
nombres différents ; on les distingue ici :
\[
q_i = \frac{(p^{i+1})!}{\bigl((p^{i})!\bigr)^{p}}, \qquad
\kappa_n = \frac{\prod_i \bigl((p^{i})!\bigr)^{a_i}}{n!}, \qquad
e_n = \frac{(pn)!}{(n!)^{p}\, p!}, \qquad
d_m = \frac{(pm)!}{(p!)^{m}\, m!},
\]
où $n = \sum a_i p^{i}$ est l'écriture de $n$ en base $p$. On note
$s_p(n) = \sum a_i$ la somme des chiffres (son « $\mathrm{chif}(n)$ »).

\emph{Sites.} Les noms changent d'un feuillet à l'autre, et c'est le piège
principal du dossier. On fixe ici, pour un $S$-schéma $X$ :

\begin{itemize}
\item $\mathrm{Cris}(X/S)$ : les épaississements $(U, U', \delta)$ d'un
ouvert $U$ de $X$, d'idéal nilpotent muni de puissances divisées $\delta$
— le site cristallin d'aujourd'hui. C'est le « $X_{\mathrm{cris}}$ » de la
page 2 et le « $\mathrm{DR}$ » des pages 27 et 28.
\item $\mathrm{Cris}^{\mathrm{st}}(X/S)$ : le sous-site des objets qui
admettent localement une rétraction $U' \to U$. C'est le
« $X_{\mathrm{conn}}$ » de la page 2 et le « $\mathrm{DR}_{\mathrm{st}}$ »
de la page 28.
\item $\mathrm{Inf}(X/S)$ : les mêmes épaississements, sans puissances
divisées — le site infinitésimal. C'est le « $X_{\mathrm{inf}}$ » de la page
2 et le « $\mathrm{Cris}$ » de la page 28.
\item $\mathrm{Inf}^{\mathrm{st}}(X/S)$ : son sous-site des objets
localement rétractables. C'est le « $X_{\mathrm{strat}}$ » de la page 2 et le
« $\mathrm{Strat}$ » de la page 28.
\end{itemize}

Cette identification est la nôtre ; elle repose sur le fait que les deux
feuillets disent la même chose des mêmes flèches.\footnote{La page 2 dit que
« $\mathrm{conn} \to \mathrm{strat}$ » est une équivalence en
caractéristique $0$ et « $\mathrm{conn} \to \mathrm{cris}$ » une équivalence
si $X/S$ est lisse ; la page 28 dit la même chose de
« $\mathrm{DR}_{\mathrm{st}} \to \mathrm{Strat}$ » et de
« $\mathrm{DR}_{\mathrm{st}} \to \mathrm{DR}$ ». Seuls les sites
$\mathrm{DR}$ et $\mathrm{DR}_{\mathrm{st}}$ sont définis sur les pages (page
27) ; $\mathrm{Strat}$ et le $\mathrm{Cris}$ de la page 28 ne le sont pas, et
leur lecture comme variantes sans puissances divisées est déduite des
énoncés qui les concernent. Qu'en 1965--66 « cristallin » ait pu désigner le
site sans puissances divisées n'a rien d'étonnant : dans les exposés de
l'IHES de 1966, « Crystals and the De Rham cohomology of schemes », les
cristaux sont d'abord définis sur le site infinitésimal, et les puissances
divisées n'arrivent qu'ensuite.} Ce qui distingue les deux topos qui portent
le nom « cristallin » dans le dossier n'est pas un détail : l'un a des
puissances divisées, l'autre non, et en caractéristique $p$ ce sont deux
théories différentes.

\subsection*{Ce que le dossier annonce sans l'établir}

Que l'anneau $W^{(\nu)}$ de la page 4 porte bien des puissances divisées ;
que la présentation de l'enveloppe des pages 31 à 34 en soit une (« à
vérifier ») ; que la cohomologie d'un objet du site stratifié soit celle de
son épaississement (page 29, « à vérifier ») ; que les foncteurs dérivés du
topos de Frobenius se calculent sur les faisceaux sous-jacents (pages 18 et
19, la feuille s'interrompt) ; qu'une opération $\gamma_p$ s'étende en une
structure de puissances divisées complète (page 40, question ouverte). Chacun
de ces points est signalé là où il se présente, avec ce que l'on en sait
aujourd'hui quand c'est sûr.

\pagerange{2}{2}
\section*{Les topos attachés à $X$ (page 2)}

Soit $X$ un schéma sur une base $S$ (sous-entendue dans les notations, dit
une note marginale). Un seul feuillet, au crayon, dessine les topos que l'on
peut attacher à $X$ et les morphismes qui les relient. Dans les notations
fixées plus haut :

\begin{tikzcd}[column sep=small, row sep=small, nodes={font=\scriptsize}]
 & & X_{\mathrm{zar}} \arrow[dll] \arrow[drr] & & \\
\mathrm{Cris}^{\mathrm{st}} \arrow[rrrr, dashed] \arrow[drr, "\text{équiv. en car. } 0"] \arrow[dd, "\text{équiv. si } X/S \text{ lisse}"'] & & & & \mathbb{P}\text{-}\mathrm{Cris}^{\mathrm{st}} \arrow[dll, dashed] \arrow[dd, "\text{équiv. si } X/p \text{ lisse}"] \\
 & & \mathrm{Inf}^{\mathrm{st}} \arrow[dd] & & \\
\mathrm{Cris} \arrow[rrrr, dashed] \arrow[drr] \arrow[ddrr] & & & & \mathbb{P}\text{-}\mathrm{Cris} \arrow[dll, dashed] \arrow[ddll] \\
 & & \mathrm{Inf} \arrow[d] & & \\
 & & X_{\mathrm{zar}} & &
\end{tikzcd}

Les flèches en tirets ne sont définies, dit une accolade, que si $X$ est de
torsion sur $\mathbf{Z}$.\footnote{Le diagramme est celui de la page, nœud
pour nœud et flèche pour flèche ; seuls les noms des sommets sont
modernisés. Les deux sommets de droite portent sur la page un $\mathbb{P}$
gras devant « conn » et « cris » ; aucun feuillet du dossier ne dit ce qu'il
désigne, et nous ne l'identifions pas. Le label de la flèche verticale de
droite, « équiv. si $X/p$ lisse », est une lecture incertaine. Un trait fin
sans pointe descend en outre, sur la page, du $X_{\mathrm{zar}}$ du haut vers
$\mathrm{Inf}$ ; il n'est pas rendu comme flèche.} Les deux équivalences de
la colonne de gauche sont vraies, et pour une raison simple que la page 28
rendra explicite : un sous-site qui contient localement tous les objets
donne le même topos. Si $X$ est lisse sur $S$, tout épaississement
$U \hookrightarrow U'$ avec $U'$ affine admet une rétraction $U' \to U$ — c'est
la définition même de la lissité formelle — de sorte que
$\mathrm{Cris}^{\mathrm{st}}$ est tout $\mathrm{Cris}$. En caractéristique
$0$, tout idéal nilpotent porte une unique structure de puissances divisées,
$\gamma_n(x) = x^{n}/n!$, de sorte que les sites avec et sans puissances
divisées sont les mêmes.

Trois remarques suivent.

Les flèches qui aboutissent à $X_{\mathrm{zar}}$ (en bas) donnent des
isomorphismes de localisation et des suites spectrales de Leray : c'est par
elles que la cohomologie cristalline se compare à la cohomologie de Zariski.

Les flèches qui partent de $X_{\mathrm{zar}}$ (en haut) donnent, sur le site
cristallin, la suite exacte
\[
0 \longrightarrow \mathcal{J}_{X/S} \longrightarrow \mathcal{O}_{X/S}
\longrightarrow \mathcal{O}_{X_0} \longrightarrow 0,
\]
où $\mathcal{O}_{X/S}$ est le faisceau structural du site, qui à un
épaississement $(U, U')$ associe $\Gamma(U', \mathcal{O}_{U'})$,
$\mathcal{O}_{X_0}$ celui qui lui associe $\Gamma(U, \mathcal{O}_U)$, et
$\mathcal{J}_{X/S}$ l'idéal à puissances divisées.\footnote{La page écrit
$\underline{\mathcal{O}}_{\mathcal{X}} \to
\underline{\mathcal{O}}_{\mathcal{X}_0}$, sans nommer le noyau autrement que
$J$, et note en marge qu'il faut « dire que
$\mathrm{Mod}(X_{\mathrm{zar}}) \simeq
\mathrm{Mod}(\underline{\mathcal{O}}_{\mathcal{X}_0})$ » — que les
modules sur $\mathcal{O}_{X_0}$ dans le topos cristallin sont les modules de
Zariski sur $X$. La remarque est rapportée comme il la fait ; elle n'est pas
démontrée sur la page.}

Enfin, pour un $\mathcal{O}_X$-module inversible $L$ sur $X_{\mathrm{zar}}$,
il définit une classe $c(L) \in H^{2}(\mathrm{Cris})$ — et de même dans
$H^{2}(\mathbb{P}\text{-}\mathrm{Cris})$, les deux étant compatibles si $X$
est de torsion — dont l'annulation est nécessaire et suffisante pour que $L$
provienne d'un module inversible sur le site cristallin, c'est-à-dire, dans
le cas relatif lisse, pour que $L$ admette une connexion intégrable. Telle
qu'on peut la reconstituer, c'est une \emph{classe d'obstruction}. La suite
exacte ci-dessus donne une suite exacte de groupes multiplicatifs
\[
1 \longrightarrow 1 + \mathcal{J}_{X/S} \longrightarrow
\mathcal{O}_{X/S}^{\times} \longrightarrow \mathcal{O}_{X_0}^{\times}
\longrightarrow 1,
\]
la classe de $L$ dans $H^{1}(\mathcal{O}_{X_0}^{\times})$ a un bord dans
$H^{2}(1 + \mathcal{J}_{X/S})$, et ce bord est nul si et seulement si $L$ se
relève en un $\mathcal{O}_{X/S}$-module inversible. Le logarithme des
puissances divisées (page 26) identifie $1 + \mathcal{J}_{X/S}$ au groupe
additif $\mathcal{J}_{X/S}$, ce qui met la classe dans
$H^{2}(\mathrm{Cris}, \mathcal{J}_{X/S})$.\footnote{Cette reconstruction est
la nôtre : la page dit seulement « on définit une classe de Chern
$c'(\underline{L})$ dans $H^{2}(X_{\mathrm{cris}})$ », sans coefficients, et
la lettre $c'$ est une lecture incertaine (peut-être $c^{1}$). Le mot
« classe de Chern » ne doit pas tromper : la première classe de Chern
cristalline d'aujourd'hui, obtenue par $d\log$, vit dans
$H^{2}(X/S, \mathcal{O}_{X/S})$ et ne s'annule pas pour cette raison-là ;
ce que la page décrit — nul si et seulement si $L$ se relève — est une
obstruction. L'usage du logarithme demande que les puissances divisées
soient localement nilpotentes, ce qui est le cas sur un site dont les
épaississements sont nilpotents.} Dans le cas lisse, un module inversible
sur le site cristallin est un module inversible muni d'une connexion
intégrable — à condition, en caractéristique $p$, que cette connexion soit
\emph{quasi-nilpotente}.\footnote{Hypothèse ajoutée. L'équivalence entre
cristaux et modules à connexion intégrable sur un $X$ lisse demande, en
caractéristique $p$, la quasi-nilpotence de la connexion ; la page dit
seulement « qu'il admette une connexion intégrable ».}

Deux questions en marge : « détermination des \emph{points} de ces topos ? »,
et le rappel que $S$ est sous-entendu partout.

\pagerange{4}{8}
\section*{L'anneau artinien universel à puissances divisées nilpotentes (pages 4 à 8)}

\subsection*{L'énoncé}

Soit $k$ un corps parfait de caractéristique $p$ et $W = W(k)$ l'anneau de
ses vecteurs de Witt. Soit $\nu \geqslant 1$. On considère la catégorie des
anneaux locaux artiniens $\Lambda$ de corps résiduel $k$ dont l'idéal maximal
$\mathfrak{r}$ est muni de puissances divisées \emph{$\nu$-nilpotentes},
c'est-à-dire telles que $\gamma_n(x) = 0$ pour $x \in \mathfrak{r}$ et
$n \geqslant p^{\nu+1}$.\footnote{La page ne définit pas « $\nu$-nilpotentes
». Les relations qu'elle impose à l'anneau universel ($\Lambda_\nu^{p} = 0$,
ci-dessous) et le calcul de la page 8 ($\Lambda^{a} = 0$ dès que $\sum a_i
p^{i} \geqslant p^{\nu+1}$) fixent le sens que nous adoptons. La condition
plus faible $\gamma_{p^{\nu}}(x)^{p} = 0$ donnerait le même objet
initial.} Il s'agit d'en construire un objet initial.

Tout tel $\Lambda$ est une $W$-algèbre, de façon unique (théorème de
structure de Cohen\footnote{La page écrit « (Cohen) », lecture incertaine ;
c'est le seul nom que la phrase puisse porter.}), et l'image de $p$ tombe dans
$\mathfrak{r}$. Les puissances divisées $\gamma_{p^{i}}(p)$ de cette image
satisfont alors, par la règle $\gamma_m(x)\gamma_n(x) = \binom{m+n}{m}
\gamma_{m+n}(x)$ itérée,
\[
\gamma_{p^{i}}(p)^{p} = q_i\, \gamma_{p^{i+1}}(p), \qquad
q_i = \frac{(p^{i+1})!}{\bigl((p^{i})!\bigr)^{p}}, \qquad
v_p(q_i) = 1 .
\]
D'où la candidate : avec des indéterminées $\Lambda_0, \ldots, \Lambda_\nu$
(qui figureront $\gamma_{p^{i}}(p)$),
\[
W^{(\nu)} = W[\Lambda_0, \ldots, \Lambda_\nu] \big/ \bigl(\Lambda_i^{p} - q_i
\Lambda_{i+1}\ (0 \leqslant i < \nu),\ \ \Lambda_\nu^{p},\ \ \Lambda_0 - p\bigr).
\]
Comme les nombres premiers $\ell \neq p$ sont inversibles dans $W$, il
suffit de connaître les $\gamma_{p^{i}}$ pour connaître tous les $\gamma_n$
(on le verra en détail à la page 31) : ces générateurs suffisent.

\textbf{Proposition.} \emph{$W^{(\nu)}$ est un anneau local artinien de corps
résiduel $k$, d'idéal maximal $J = (\Lambda_0, \ldots, \Lambda_\nu)$, et
c'est l'objet initial de la catégorie.}

Qu'il soit artinien est immédiat : $p^{p^{\nu+1}} = \Lambda_0^{p^{\nu+1}}$
est un multiple de $\Lambda_\nu^{p} = 0$, donc $W^{(\nu)}$ est une algèbre
finie sur $W/p^{N}$, engendrée par des nilpotents. Qu'il soit initial
\emph{parmi les anneaux qui satisfont ses relations} l'est aussi. Le point
qui demande un argument, et que la page affirme sans le donner, est que
$J$ porte bien des puissances divisées $\nu$-nilpotentes.\footnote{La page
place l'énoncé d'universalité entre crochets, sans démonstration. Il est
vrai : $W^{(\nu)}$ est le quotient de l'enveloppe à puissances divisées de
l'idéal $pW$ (celle des pages 23 et 31) par l'idéal engendré par les
$\gamma_n(p)$, $n \geqslant p^{\nu+1}$, qui est stable par les $\gamma_m$, et
l'on vérifie à partir de la formule de la page 31 que cet idéal est engendré
par les $\Lambda_j = \gamma_{p^{j}}(p)$, $j \geqslant \nu + 1$ ; le quotient
est l'anneau de la page, la relation $\Lambda_\nu^{p} = q_\nu
\Lambda_{\nu+1} = 0$ en tenant lieu. Que
l'enveloppe d'un idéal principal se calcule par changement de base à partir
de celle de $(X) \subset \mathbf{Z}_{(p)}[X]$ tient à ce qu'une structure de
puissances divisées sur un idéal principal s'étend à toute algèbre.}

\subsection*{La forme normale}

Pour calculer $W^{(\nu)}$ on part de l'anneau plus simple
\[
A^{(\nu)} = W[\Lambda_0, \ldots, \Lambda_\nu] \big/ \bigl(\Lambda_i^{p} - q_i
\Lambda_{i+1}\ (0 \leqslant i < \nu),\ \ \Lambda_\nu^{p}\bigr),
\]
sans la relation $\Lambda_0 = p$. C'est un $W$-module libre de base les
monômes $\Lambda^{a} = \Lambda_0^{a_0}\cdots\Lambda_\nu^{a_\nu}$ à
exposants $0 \leqslant a_i \leqslant p-1$. On pose
\[
\lambda_i = \frac{p^{p^{i}}}{(p^{i})!} \in W \qquad \text{et} \qquad
\Lambda_i = \lambda_i + \Theta_i ,
\]
où $\lambda_i$ est la puissance divisée \emph{usuelle} $\gamma_{p^{i}}(p)$ de
$p$ dans $W$, et où $\Theta_i$ mesure donc l'écart entre les puissances
divisées universelles et les usuelles. Les monômes
$\Theta^{a}$, $0 \leqslant a_i \leqslant p-1$, forment encore une base, et
$W^{(\nu)} = A^{(\nu)}/I$ où $I$ est l'idéal engendré par $\Theta_0$.

\textbf{Lemme (page 5).} \emph{Pour que $\sum_a w_a\, \Theta^{a}$ soit dans
$I$, il faut et il suffit que}
\begin{enumerate}
\item[a)] \emph{pour tout $a \neq 0$, si $a_i$ est la première composante non
nulle de $a$, on ait $w_a \in p^{i} W$ ;}
\item[b)] \emph{le coefficient $w_0$ de $1$ soit dans $p^{c(\nu)}W$, où}
\[
c(\nu) = \nu + p\, v_p(\lambda_\nu) = \nu + p\Bigl[p^{\nu} -
\bigl(p^{\nu-1} + \cdots + p + 1\bigr)\Bigr].
\]
\end{enumerate}

\emph{Suffisance.} Il suffit de voir que $p^{i}\Theta_i \in I$ pour
$0 \leqslant i \leqslant \nu$ et que $p^{c(\nu)} \in I$. La première relation
se prouve par récurrence sur $i$ ; pour $i = 0$ elle dit
$\Theta_0 \in I$.\footnote{La page écrit « si $i = \nu$ » pour le début de
la récurrence ; c'est $i = 0$.} Si $p^{i}\Theta_i \in I$ avec $i < \nu$, on
développe $(\lambda_i + \Theta_i)^{p} = q_i(\lambda_{i+1} + \Theta_{i+1})$ ;
comme $\lambda_i^{p} = q_i \lambda_{i+1}$, il reste
\[
\sum_{\alpha=1}^{p} \binom{p}{\alpha}\, \Theta_i^{\alpha}\, \lambda_i^{p-\alpha}
= q_i\, \Theta_{i+1} .
\]
Multiplié par $p^{i}$, le membre de gauche est dans $I$, donc
$p^{i} q_i \Theta_{i+1} \in I$, et comme $v_p(q_i) = 1$, $p^{i+1}\Theta_{i+1}
\in I$.\footnote{La page écrit « $= 0$ » pour « $\in I$ » : elle calcule
modulo $I$. Un terme barré d'un trait diagonal dans le développement se lit
$\Theta_{i+1}$ où l'on attend $\lambda_{i+1}$ ; le calcul ci-dessus est le
calcul correct.} Pour la seconde, on écrit $(\lambda_\nu + \Theta_\nu)^{p} =
0$, on multiplie par $p^{\nu}$, et l'on obtient $p^{\nu}\lambda_\nu^{p} \in
I$, de valuation $\nu + p\, v_p(\lambda_\nu) = c(\nu)$, puisque
$v_p(\lambda_\nu) = p^{\nu} - v_p\bigl((p^{\nu})!\bigr) = p^{\nu} -
\frac{p^{\nu}-1}{p-1}$.

\emph{Nécessité.} L'ensemble $I'$ des éléments qui satisfont a) et b) est
un sous-$W$-module contenu dans $I$ et contenant $\Theta_0$ ; il suffit
de voir que c'est un idéal, c'est-à-dire qu'il est stable par multiplication
par les $\Theta_i$. La page laisse la vérification au lecteur, avec un point
d'exclamation en marge.\footnote{Nous ne l'avons pas faite à la main. Nous
avons vérifié par le calcul — en exprimant la multiplication par
$\Theta_0$ dans la base des $\Theta^{a}$ et en comparant les réseaux —
que $I$ est exactement le réseau diagonal décrit par a) et b), pour
$(p, \nu) = (2,1), (2,2), (2,3), (3,1), (3,2), (5,1)$. Pour $p = 2$, $\nu = 1$,
on trouve $W^{(1)} = W[\Theta_1]/(2\Theta_1, \Theta_1^{2} + 4)$, de longueur
$4$, avec $c(1) = 3$.}

\textbf{Corollaire (page 6).} \emph{Soit $E_i$ un système de représentants de
$W/p^{i}W$. Tout élément de $W^{(\nu)}$ s'écrit de façon unique}
\[
\sum w_{a_1 \ldots a_\nu}\, \Theta_1^{a_1}\cdots\Theta_\nu^{a_\nu},
\qquad 0 \leqslant a_i \leqslant p-1,
\]
\emph{avec $w_a \in E_i$ si $a_i$ est la première composante non nulle de
$a$, et $w_0 \in E_{c(\nu)}$.} En particulier $W^{(\nu)}$ est de longueur
$c(\nu) + \sum_{i=1}^{\nu} i\,(p-1)\,p^{\nu-i}$ sur $W$.\footnote{Ce compte
de longueur est une conséquence immédiate du corollaire ; il n'est pas sur la
page.}

\textbf{Lemme (page 7).} Pour $1 \leqslant r \leqslant \nu$, soit
$W^{(\nu, r)}$ le quotient de $W^{(\nu)}$ par $\Theta_1, \ldots, \Theta_r$ :
c'est l'anneau universel où les puissances divisées de $p$ coïncident avec
les usuelles jusqu'à l'ordre $p^{r}$. Avec $I_{\nu,r} = (\Theta_0, \ldots,
\Theta_r)$, de sorte que $A^{(\nu)}/I_{\nu,r} = W^{(\nu, r)}$ : \emph{pour que
$\sum w_a \Theta^{a}$ soit dans $I_{\nu,r}$, il faut et il suffit que}
\begin{enumerate}
\item[a)] \emph{pour tout $a \neq 0$ dont la première composante non nulle
$a_i$ a un indice $i > r$, on ait $w_a \in p^{i-r}W$} (aucune condition si
$i \leqslant r$) ;
\item[b)] \emph{$w_0 \in p^{c(\nu)-r}W$.}
\end{enumerate}
D'où la forme normale
$\sum w_a\, \Theta_{r+1}^{a_{r+1}}\cdots\Theta_\nu^{a_\nu}$, avec $w_a \in
E_{i}$ si $a_{r+i}$ est la première composante non nulle, et $w_0 \in
E_{c(\nu)-r}$.\footnote{La parenthèse « aucune condition si $i \leqslant r$ »
est implicite sur la page, qui n'énonce a) que pour $i > r$. Vérifié par le
même calcul que le lemme précédent, pour tous les $r$ des cas cités. Pour
$r = \nu$ on trouve $W^{(\nu,\nu)} = W/p^{p\,v_p(\lambda_\nu)}$ : si les
puissances divisées de $p$ sont les usuelles, la seule relation est
$\lambda_\nu^{p} = 0$.}

\subsection*{Le calcul indépendant de la page 8}

La page 8 reprend l'anneau $A^{(\nu)}$ pour lui-même et montre que
\[
\Lambda_0^{a_0}\cdots\Lambda_\nu^{a_\nu} = 0 \qquad \text{dès que} \qquad
\sum a_i p^{i} \geqslant p^{\nu+1} \qquad (a_i \geqslant 0).
\]
C'est ce qu'on attend si $\Lambda^{a}$ figure, à une constante près,
$\gamma_{\sum a_i p^{i}}(\Lambda_0)$. La preuve est une retenue en base $p$.
\emph{Lemme} : $\Lambda^{a} = c\,\Lambda^{b}$ avec $c \in W$,
$0 \leqslant b_i \leqslant p-1$ pour $i < \nu$, et $\sum a_i p^{i} = \sum b_i
p^{i}$. Si $a_i \geqslant p$ pour un $i < \nu$, on écrit $a_i = p\,q + a'_i$ ;
alors $\Lambda_i^{a_i} = (q_i\Lambda_{i+1})^{q}\Lambda_i^{a'_i}$, ce qui
remplace $(a_i, a_{i+1})$ par $(a'_i, a_{i+1} + q)$ sans changer
$\sum a_i p^{i}$ ; on recommence jusqu'à ce que tous les exposants d'indice
$< \nu$ soient $< p$. Si alors $\sum b_i p^{i} \geqslant p^{\nu+1}$, c'est que
$b_\nu \geqslant p$, et $\Lambda_\nu^{b_\nu}$ est un multiple de
$\Lambda_\nu^{p} = 0$.\footnote{Deux lignes biffées ouvrent la démonstration,
et une note encadrée en marge (« si tous les $a_i$ sont $\leqslant p - 1$,
c'est ok ») en donne le cas trivial. Les exposants de la ligne principale
sont écrits petit ; ils sont restitués d'après $a_i = pq_i + a'_i$.}

\pagerange{10}{10}
\section*{La formule de Legendre (page 10)}

Un demi-feuillet en travers calcule $v_p(n!)$ à partir des $v_p\bigl((p^{k})!
\bigr)$ :
\[
v_p\bigl((p^{k})!\bigr) = 1 + p + \cdots + p^{k-1} = \frac{p^{k}-1}{p-1},
\]
puis, pour $n = \sum a_i p^{i}$,
\[
v_p(n!) = \sum_i a_i\, v_p\bigl((p^{i})!\bigr), \qquad \text{d'où} \qquad
\boxed{\ v_p(n!) = \frac{n - s_p(n)}{p-1}\ }
\]
— la formule de Legendre, encadrée deux fois sur la page.\footnote{La page
écrit $v_p(p)$, $v_p(p^{2})$, \ldots\ pour $v_p(p!)$, $v_p\bigl((p^{2})!\bigr)$,
\ldots\ La première ligne, qui exprime $v_p(n!)$ en termes d'une coupure en
$\mu$, n'est pas reliée par la page à la suite et n'est pas reprise ici. Le
nom de Legendre est en marge de la page 46, pour la conséquence
$\sup_n v_p(n!)/n = 1/(p-1)$.}

\pagerange{12}{12}
\section*{Quand $\pi^{n}/n!$ est-il entier, et quand tend-il vers zéro ? (page 12)}

Soit $V$ un anneau de valuation, $v_p$ sa valuation normalisée par
$v_p(p) = 1$, et $\pi \in V$. D'après Legendre,
\[
(p-1)\, v_p\Bigl(\frac{\pi^{n}}{n!}\Bigr) = n\bigl[(p-1)v_p(\pi) - 1\bigr]
+ s_p(n).
\]
Il en résulte :
\[
v_p\Bigl(\frac{\pi^{n}}{n!}\Bigr) \geqslant 0 \ \text{pour tout } n
\iff v_p\Bigl(\frac{\pi^{n}}{n!}\Bigr) > 0 \ \text{pour tout } n \geqslant 1
\]
\[
\iff v_p(\pi) \geqslant \frac{1}{p-1}
\iff \pi^{p-1} \in pV ,
\]
\[
v_p\Bigl(\frac{\pi^{n}}{n!}\Bigr) \longrightarrow +\infty
\iff v_p(\pi) > \frac{1}{p-1} .
\]
Si $v_p(\pi) = 1/(p-1)$, la quantité vaut $s_p(n)/(p-1)$, qui reprend la
valeur $1/(p-1)$ pour tous les $n = p^{k}$ et ne tend donc pas vers
l'infini ; si $v_p(\pi) < 1/(p-1)$, le choix $n = p^{k}$ la rend négative
pour $k$ grand. Lorsque $p = u\,\pi^{r}$ avec $u$ inversible — par exemple
si $V$ est de valuation discrète et $\pi$ une uniformisante, $r$ étant
alors l'indice de ramification —, la seconde condition s'écrit
$r < p - 1$, et la première $r \leqslant p-1$ : c'est l'inégalité de la page
46.\footnote{La page établit $s_p(n)/n \to 0$ en majorant
$s_p(n) \leqslant (p-1)\nu$ pour $p^{\nu} \leqslant n < p^{\nu+1}$ ; la
borne exacte est $(p-1)(\nu+1)$, ce qui ne change rien à la conclusion. Elle
écrit aussi « $\mathrm{chif}(n) \leqslant 1$ » pour $\mathrm{chif}(n)/n
\leqslant 1$. Elle ne donne que la seconde équivalence sous la forme « $p =
\text{unit}\cdot\pi^{r}$, avec $r < p-1$ » ; la forme $r \leqslant p-1$ de la
première est la nôtre.}

C'est le critère qui, aujourd'hui, distingue les puissances divisées
\emph{qui existent} de celles qui sont \emph{topologiquement nilpotentes} sur
l'idéal maximal d'un anneau de valuation discrète : $e \leqslant p-1$ pour
les premières, $e < p-1$ pour les secondes.

\pagerange{15}{16}
\section*{Frobenius : le topos $\mathcal{Q}(\alpha, \beta)$ (pages 15 à 19)}

Une chemise de sa main porte « Frobenius » ; le premier feuillet s'intitule
« Général nonsens sur Frobenius ».

\subsection*{1. Faisceaux munis d'une flèche de Frobenius}

Soit un carré commutatif de topos
\begin{tikzcd}
X \arrow[r, "F_X"] \arrow[d, "g"'] & X \arrow[d, "g"] \\
Y \arrow[r, "F_Y"] & Y
\end{tikzcd}
(on pense aux Frobenius de $X$ et de $Y$, mais rien ne l'exige). Pour un
faisceau $E$ sur $X$ on dispose du morphisme de changement de base
\[
(1.1) \qquad c_E : F_Y^{*}\, g_*(E) \longrightarrow g_*\, F_X^{*}(E),
\]
défini pour tout carré commutatif par adjonction. Si $E$ est muni d'une
flèche $F_E : F_X^{*}E \to E$, on en tire
\[
(1.3) \qquad g_*(F_E) \circ c_E : F_Y^{*}\bigl(g_*E\bigr) \longrightarrow
g_*E ,
\]
ce qui munit $g_*E$ d'une flèche de même nature ; si $E$ est un faisceau de
groupes, d'anneaux, de modules, et $F_E$ un morphisme de tels, il en va de
même de (1.3), puisque (1.1) l'est. Dans l'autre sens, pour un faisceau $G$
sur $Y$ muni de $F_G : F_Y^{*}G \to G$, l'isomorphisme
$c^{G} : F_X^{*}g^{*}G \xrightarrow{\sim} g^{*}F_Y^{*}G$ donne
\[
(1.6) \qquad g^{*}(F_G) \circ c^{G} : F_X^{*}\bigl(g^{*}G\bigr)
\longrightarrow g^{*}G .
\]
Autrement dit, les foncteurs $g_*$ et $g^{*}$ se relèvent aux faisceaux
munis d'une flèche de Frobenius. C'est le langage des $F$-cristaux, avant
la lettre.

\subsection*{2. Le topos $\mathcal{Q}(\alpha, \beta)$}

Soient $\alpha, \beta : X' \rightrightarrows X$ deux morphismes de topos. Soit
$\mathcal{Q}(\alpha, \beta)$ la catégorie des couples $(E, F_E)$ formés d'un
faisceau $E$ sur $X$ et d'une flèche $F_E : \alpha^{*}E \to \beta^{*}E$
— variante : un \emph{isomorphisme} —, et
$\omega^{*} : \mathcal{Q}(\alpha,\beta) \to X$ le foncteur d'oubli.

\textbf{Proposition.} \emph{$\mathcal{Q}(\alpha, \beta)$ est un topos, et
$\omega^{*}$ est le foncteur image inverse d'un morphisme de topos
$\omega : X \to \mathcal{Q}(\alpha, \beta)$.}

Les flèches $F_E$, naturelles en $(E, F_E)$, forment une transformation
$F : \omega\beta \to \omega\alpha$ entre morphismes
$X' \to \mathcal{Q}(\alpha,\beta)$, et le couple $(\omega, F)$ est
universel : pour tout topos $Y$, le foncteur
\[
g \longmapsto (g \circ \omega,\ g * F) \;:\;
\underline{\mathrm{Homtop}}\bigl(\mathcal{Q}(\alpha, \beta), Y\bigr)
\longrightarrow \underline{\mathrm{Homtop}}(X, \alpha, \beta ; Y)
\]
est une équivalence, où le second membre est la catégorie des couples
$(\varphi, f)$ d'un morphisme de topos $\varphi : X \to Y$ et d'une flèche
$f : \varphi\beta \to \varphi\alpha$.\footnote{Convention : une flèche
$\varphi \to \psi$ entre morphismes de topos est ici une flèche
$\varphi_* \to \psi_*$, c'est-à-dire $\psi^{*} \to \varphi^{*}$ ; c'est avec
elle que $F_E : \alpha^{*}\omega^{*} \to \beta^{*}\omega^{*}$ est une flèche
$\omega\beta \to \omega\alpha$, comme l'écrit la page. La page place $F$ dans
$\underline{\mathrm{Homtop}}(X, \mathcal{Q}(\alpha,\beta))$ ; c'est
$\underline{\mathrm{Homtop}}(X', \mathcal{Q}(\alpha,\beta))$, les composés
$\omega\alpha$, $\omega\beta$ partant de $X'$. Un crochet serré en marge
(« qui se \ldots\ à l'aide de $\underline{\mathrm{Hom}}(X, Y)
\rightrightarrows \underline{\mathrm{Homtop}}(X', Y)$ ») n'est pas lu en
entier.}

Dans le langage d'aujourd'hui, $\mathcal{Q}(\alpha, \beta)$ est
l'\emph{inserteur} des foncteurs $\alpha^{*}, \beta^{*} : X \to X'$ dans la
2-catégorie des catégories (dans la variante inversible, l'iso-inserteur),
et, comme topos, le \emph{co-inserteur} — une colimite lâche — des morphismes
$\alpha, \beta$ dans la 2-catégorie des topos.\footnote{Ces noms sont
postérieurs : les limites pondérées de ce type (inserteurs, iso-inserteurs)
ont été nommées et étudiées par l'école australienne de théorie des
catégories (Kelly, Street) dans les années 1970 et 1980. La page n'emploie que
la description universelle, qui est la même chose.} L'exemple le plus petit
le montre : si $X = X' = \mathbf{Ens}$ et $\alpha = \beta = \mathrm{id}$, un
objet est un ensemble muni d'un endomorphisme, et $\mathcal{Q}$ est le topos
classifiant du monoïde $\mathbf{N}$ (du groupe $\mathbf{Z}$ dans la variante
inversible).\footnote{Exemple de l'édition.}

\pagerange{17}{17}
\subsection*{Pourquoi c'est un topos}

On vérifie les conditions de Giraud : les limites finies et les colimites de
$\mathcal{Q}$ se calculent sur les faisceaux sous-jacents, puisque
$\alpha^{*}$ et $\beta^{*}$ commutent aux unes et aux autres, et il reste à
exhiber une petite famille génératrice. Dans la variante inversible,
$\mathcal{Q}$ s'interprète comme la catégorie des sections cartésiennes d'un
topos fibré au-dessus de la catégorie à deux objets $x', x$ et deux flèches
$\alpha, \beta : x' \rightrightarrows x$, de fibres $X'$ en $x'$ et $X$ en $x$,
les changements de base étant $\alpha^{*}$ et $\beta^{*}$ ; dans la variante
non inversible, une catégorie fibrée analogue, avec un sommet de plus, fait
l'affaire, sous réserve de vérification.\footnote{La page renvoie à « SGA 4
I 9.25 » pour le premier cas et à « SGA 4 I 9.21 et suivants » pour le
second ; nous reproduisons ces renvois sans les avoir vérifiés. Le diagramme
marginal du second cas ($x' \to x$ par $u$, $x' \to x'_1$ par $f$,
$x'_1 \to x$ par $v$) est encadré et accompagné de « à vérifier ». Plusieurs
mots de ce paragraphe sont illisibles ; l'argument par Giraud est restitué
d'après ce qui se lit.}

Le cas qui motive tout est $X' = X$, $\beta = \mathrm{id}$ : les objets sont
alors les $(E, F_E)$ avec $F_E : \alpha^{*}E \to E$, exactement ceux du n° 1
quand $\alpha$ est un Frobenius.

\pagerange{17}{19}
\subsection*{Fonctorialité, et la question des foncteurs dérivés}

Le topos $\mathcal{Q}(\alpha,\beta)$ dépend fonctoriellement de la paire
$\alpha, \beta$. Soit un diagramme de topos, commutatif à isomorphisme près,
\begin{tikzcd}
X' \arrow[r, "{\alpha,\ \beta}"] \arrow[d, "u_1"'] & X \arrow[d, "u_0"] \\
Y' \arrow[r, "{\lambda,\ \mu}"] & Y
\end{tikzcd}
c'est-à-dire avec $u_0\alpha \simeq \lambda u_1$ et $u_0\beta \simeq \mu
u_1$.\footnote{La page trace des flèches doubles, $\alpha, \beta$ en haut et
$\lambda, \mu$ en bas ; le diagramme n'en rend qu'une par ligne, qui porte
les deux noms.} La propriété universelle fournit
un morphisme $v = \mathcal{Q}(u) : \mathcal{Q}(\alpha,\beta) \to
\mathcal{Q}(\lambda,\mu)$, avec les transitivités attendues, tel que
$v\,\omega_X \simeq \omega_Y\, u_0$. Dans le cas $X' = X$, $Y' = Y$,
$\beta = \mathrm{id}_X$, $\mu = \mathrm{id}_Y$, $\alpha = F_X$, $\lambda =
F_Y$, on retrouve le carré du n° 1, et $v_*$, $v^{*}$ sont donnés par les
formules (1.3) et (1.6) : $v_*(E, F_E) = (u_{0*}E,\ u_{0*}(F_E) \circ c_E)$,
$v^{*}(G, F_G) = (u_0^{*}G,\ u_0^{*}(F_G)\circ c^{G})$.

Pour un objet $\mathbb{E} = (E, F_E)$ de $\mathcal{Q}(\alpha, \beta)$ —
disons un faisceau abélien —, on dispose des morphismes de changement de base
\[
\omega_Y^{*}\, R^{i}v_*(\mathbb{E}) \longrightarrow R^{i}u_{0*}\bigl(
\omega_X^{*}\mathbb{E}\bigr) = R^{i}u_{0*}(E),
\]
et l'on aimerait savoir que ce sont des isomorphismes, ce qui est le cas pour
$i = 0$ par la formule ci-dessus.\footnote{La page écrit $\omega_X^{*}$ à
gauche ; c'est $\omega_Y^{*}$, puisque $R^{i}v_*(\mathbb{E})$ est un objet
de $\mathcal{Q}(\lambda,\mu)$.} Comme $\omega_Y^{*}$ est exact, il suffit
que $\omega_X^{*}$ envoie les injectifs de $\mathcal{Q}(\alpha,\beta)$ sur
des faisceaux acycliques pour $u_{0*}$. Il affirme qu'il suffit pour cela que
$\alpha^{*}$ et $\beta^{*}$ commutent aux produits infinis — que $\alpha_!$ et
$\beta_!$ existent —, auquel cas $\omega_X^{*}$ commute aussi aux produits et
$\omega_{X!}$ existe ; la feuille s'interrompt là, au milieu de la ligne, et
l'argument n'est pas écrit.\footnote{Les lignes interlinéaires de la page 18
qui explicitaient $v_*$ et $v^{*}$ « comme au n° 1 » ne se lisent pas ; les
formules données ci-dessus sont celles que (1.3) et (1.6) imposent. La page
19 ne porte qu'une ligne, qui clôt la page 18, et sa fin est perdue dans la
marge.}

\pagerange{21}{22}
\section*{Pseudo-anneaux à puissances divisées (pages 21 à 26)}

Les pages 21 à 30 sont une rédaction numérotée, §§ 1 à 4 de sa main.

\subsection*{§ 1. Les axiomes}

Un \emph{pseudo-anneau} est un groupe abélien muni d'une multiplication
associative, bilinéaire — et ici commutative — mais sans unité : c'est le
statut naturel d'un idéal. Des puissances divisées sur un pseudo-anneau $J$
sont des applications $\gamma_n : J \to J$, $n \geqslant 1$, telles que
\begin{enumerate}
\item[1.1.] $\gamma_1(x) = x$ ;
\item[1.2.] $\gamma_n(x+y) = \gamma_n(x) + \sum_{p+q=n,\ p, q \geqslant 1}
\gamma_p(x)\gamma_q(y) + \gamma_n(y)$ ;
\item[1.3.] $\gamma_n(xy) = x^{n}\gamma_n(y)$ ;
\item[1.4.] $\gamma_m(\gamma_n(x)) = \dfrac{(mn)!}{(n!)^{m}\, m!}\,
\gamma_{mn}(x)$ ;
\item[1.5.] $\gamma_m(x)\gamma_n(x) = \dbinom{m+n}{m}\gamma_{m+n}(x)$.
\end{enumerate}
Si $J$ est un idéal d'un anneau $A$, on pose $\gamma_0 = 1$ et 1.2 s'écrit
$\gamma_n(x+y) = \sum_{p+q=n}\gamma_p(x)\gamma_q(y)$. En itérant 1.5,
\[
\text{(1.5 bis)} \qquad \gamma_{p_1}(x)\cdots\gamma_{p_r}(x) =
\frac{(p_1 + \cdots + p_r)!}{p_1!\cdots p_r!}\,\gamma_{p_1+\cdots+p_r}(x),
\qquad \text{(1.5 ter)} \qquad n!\,\gamma_n(x) = x^{n}.
\]
Si $J$ est de plus une algèbre sur un anneau $A$, on demande
\begin{enumerate}
\item[1.6.] $\gamma_n(\lambda x) = \lambda^{n}\gamma_n(x)$ pour $\lambda \in
A$, $x \in J$,
\end{enumerate}
condition automatique si $A = \mathbf{Z}$.\footnote{La note marginale dit
« c'est tjs vérifié si $A = \mathbf{Z}$ » ; c'est vrai, par récurrence à
partir de 1.2 et 1.5 ($\gamma_n(2x) = 2^{n}\gamma_n(x)$, etc.). La page
écrit « $x \in A$ » biffé dans 1.6 ; c'est $x \in J$.} Des puissances
divisées sur un idéal $J$ de $A$ sont alors la même chose que des puissances
divisées sur la $A$-algèbre sous-jacente à $J$.

Ce sont, à la présentation près, les axiomes de Cartan (1954) et de Roby
(1963), que Berthelot reprendra dans sa thèse.\footnote{La page ne cite
personne.}

\subsection*{§ 2. L'algèbre $\Gamma(M)$}

Pour un $A$-module $M$, l'algèbre $\Gamma_A(M) = \bigoplus_{n \geqslant 0}
\Gamma^{n}(M)$ est munie des applications $x \mapsto x^{[n]}$ de $M$ dans
$\Gamma^{n}(M)$. Il les prolonge en des $\gamma_n : \Gamma^{+}(M) \to
\Gamma^{+}(M)$, en notant que sur $\Gamma^{p}(M)$ il suffit de se donner une
application polynomiale homogène de degré $np$, à savoir
\[
x \longmapsto \bigl(x^{[p]}\bigr)^{[n]} = \frac{(pn)!}{(p!)^{n}\, n!}\,
x^{[pn]} ,
\]
et l'on vérifie 1.1 à 1.6. \emph{Ainsi $\Gamma^{+}_A(M)$ est l'algèbre à
puissances divisées universelle munie d'une application $A$-linéaire de
$M$.}\footnote{C'est le théorème de N. Roby (« Lois polynômes et lois
formelles en théorie des modules », Ann. Éc. Norm. Sup. 1963, et « Les
algèbres à puissances divisées », Bull. Sc. Math. 1965), que la page
utilise sans le nommer ; le passage par les applications polynomiales est
exactement celui de Roby. Le degré « $np$ » est une lecture incertaine ;
c'est le seul possible.}

\pagerange{23}{25}
\subsection*{§ 3. L'enveloppe à puissances divisées d'un idéal}

Soit $I$ un idéal de $A$. On cherche un anneau $B$ muni d'un idéal $J$ à
puissances divisées et un homomorphisme $u : A \to B$ tel que $u(I) \subset
J$, \emph{universel} pour ces données. C'est ce qu'on appelle aujourd'hui
l'\emph{enveloppe à puissances divisées} $D_A(I)$ de
$I$.\footnote{La page la note $\Gamma^{*}(A, I)$ à la page 23 et
$\Gamma(A, J)$ à la page 24. Le nom d'« enveloppe à puissances divisées » et
la notation $D_A(I)$ sont ceux de Berthelot (1974) et de Berthelot--Ogus
(1978).} La construction proposée :
\begin{enumerate}
\item on forme $\Gamma^{+}_A(I)$, l'idéal $I$ étant vu comme simple
$A$-module ;
\item on le divise par l'idéal engendré par les $(xy)^{[n]} - x^{n} *
y^{[n]}$, pour $x, y \in I$, où $*$ est le produit de $\Gamma_A(I)$ et $xy$
le produit de $I$ : il constate que les $\gamma_n$ passent au quotient
$J'$, qui est une $A$-algèbre à puissances divisées, et l'application
canonique $I \to J'$, $x \mapsto x'$, est compatible aux produits ;
\item on recolle $A$ et $J'$ le long de $I$ : $B$ est le quotient de
l'anneau $A \oplus J'$ (produit $(a,\xi)(b,\eta) = (ab, a\eta + b\xi +
\xi\eta)$) par l'image de $x \mapsto (x, -x')$, et $J$ est l'image de
$J'$.\footnote{La lettre dans « $\Gamma^{+}_A(\ )$ » est surchargée, $I$ ou $J$ ;
la construction demande $I$. La page écrit « $(x.y)^{(n)} - x^{n} *
y^{(n)}$ », le point désignant « la multiplication de $J$ » — c'est-à-dire,
pour ce que fait la construction, celle de l'idéal $I$. Que l'idéal du point
2 soit stable par les $\gamma_n$ est affirmé (« on constate »), non
démontré ; c'est vrai, et c'est l'argument de Berthelot--Ogus, Théorème
3.19.}
\end{enumerate}
Si $A$ est une $\Lambda$-algèbre, $D_A(I)$ l'est aussi.

\emph{Cas particulier.} Pour $A = \mathrm{Sym}_\Lambda(M)$ et $I =
\mathrm{Sym}^{+}_\Lambda(M)$, le problème universel est celui d'envoyer $M$
dans une $\Lambda$-algèbre à puissances divisées, de sorte que
\[
D_{\mathrm{Sym}_\Lambda(M)}\bigl(\mathrm{Sym}^{+}_\Lambda(M)\bigr) \simeq
\Gamma_\Lambda(M).
\]

\emph{Quotients.} Soit $(A, J, \gamma)$ un anneau à puissances divisées et
$(x_i)$ une famille d'éléments de $J$. L'idéal $K$ engendré par les
$\gamma_n(x_i)$, $n \geqslant 1$, est le plus petit idéal contenant les $x_i$
et stable par les $\gamma_n$ ; les $\gamma_n$ passent au quotient
$(A/K, J/K)$, qui représente le foncteur des morphismes d'anneaux à
puissances divisées $(A, J) \to (A', J')$ qui annulent les $x_i$.

\pagerange{26}{26}
\subsection*{Puissances divisées nilpotentes ; logarithme et exponentielle}

Les puissances divisées sont dites \emph{nilpotentes} s'il existe $n$ tel
que $\gamma_i(x) = 0$ pour $i \geqslant n$ et tout $x \in J$, \emph{localement
nilpotentes} si un tel $n$ existe pour chaque $x$.\footnote{La page porte
trois formulations superposées et en partie biffées ; elle parle aussi de
« nilidéal », et note que si $J$ est engendré par $N$ éléments, des
puissances divisées nilpotentes entraînent $J^{N'} = 0$ pour un $N'$
convenable. Ce que le logarithme et l'exponentielle demandent est la
nilpotence locale, et c'est la définition que la page retient en dernier
(« $\forall x$, $\exists n$ tel que $x^{(i)} = 0$ si $i > n$ »). Le nom
moderne prête à confusion : chez Berthelot, « PD-nilpotent » désigne une
autre condition, $J^{[N]} = 0$ pour la filtration par les puissances divisées.}
Dans ce cas $J$ est un nilidéal ($x^{n} = n!\gamma_n(x)$), $1 + J$ est
formé d'inversibles, et l'on a des isomorphismes de groupes réciproques
\[
\exp : J \xrightarrow{\ \sim\ } 1 + J, \quad \exp x = \sum_{n \geqslant 0}
\gamma_n(x) ;
\]
\[
\log : 1 + J \xrightarrow{\ \sim\ } J, \quad \log(1+x) = \sum_{n \geqslant 1}
(-1)^{n-1}(n-1)!\,\gamma_n(x) ,
\]
toutes les sommes étant finies.\footnote{La page écrit la série du
logarithme à partir de $n \geqslant 0$ et avec le signe $(-1)^{n}$ ; la
borne est $n \geqslant 1$ et le signe $(-1)^{n-1}$, comme le veut
$x^{n}/n = (n-1)!\,x^{n}/n!$.} C'est ce logarithme qui sert, à la page 2, à
passer du groupe multiplicatif $1 + \mathcal{J}$ au groupe additif
$\mathcal{J}$.

\pagerange{27}{28}
\section*{Le site de De Rham (pages 27 à 30)}

\subsection*{§ 4. Le site et ses variantes}

Soit $f : X \to S$. Les objets du site --- que la page appelle « site de De
Rham », et qu'on note ici $\mathrm{Cris}(X/S)$ --- sont les quadruplets
$(U, U', i, \delta)$ : $U$ un ouvert de $X$, $U'$ un $S$-schéma,
$i : U \hookrightarrow U'$ une $S$-immersion nilpotente, $\delta$ une
structure de puissances divisées sur l'idéal $\mathcal{J} =
\mathrm{Ker}(\mathcal{O}_{U'} \to \mathcal{O}_U)$. Les morphismes sont les
évidents, la topologie celle de Zariski.\footnote{La page dit les
puissances divisées « topologiquement nilpotentes (ou quasi-nilpotentes) »,
lecture par le sens d'un mot abrégé ; l'immersion étant nilpotente, la
question ne se pose que pour l'idéal $\mathcal{J}$, qui l'est. Un premier
quadruplet $(U, \mathcal{A}, \varphi, \pi^{*})$ est biffé.} À la condition
de compatibilité près avec des puissances divisées données sur la base, que
la page ne mentionne pas, c'est le site cristallin de Berthelot, dans sa
version à épaississements nilpotents.

Le site \emph{stratifié} $\mathrm{Cris}^{\mathrm{st}}(X/S)$ est formé des
objets au-dessus desquels il existe localement une rétraction $U' \to U$ de
$i$, cette rétraction ne faisant pas partie de la donnée. On a les topos
correspondants, et le carré de morphismes de topos
\begin{tikzcd}
X_{\mathrm{zar}} \arrow[r, dashed] & \mathrm{Cris}^{\mathrm{st}} \arrow[r] \arrow[d] & \mathrm{Inf}^{\mathrm{st}} \arrow[d] \\
& \mathrm{Cris} \arrow[r] & \mathrm{Inf}
\end{tikzcd}
dont il affirme :\footnote{Sur la page, les sommets sont
$\widetilde{\mathrm{Zar}}(X)$, $\widetilde{\mathrm{DR}}_{\mathrm{st}}(X)$,
$\widetilde{\mathrm{Strat}}(X)$, $\widetilde{\mathrm{DR}}(X)$,
$\widetilde{\mathrm{Cris}}(X)$ ; voir les conventions pour la
correspondance des noms. La flèche $X_{\mathrm{zar}} \to
\mathrm{Cris}^{\mathrm{st}}$ est en pointillé sur la page.}
\begin{itemize}
\item si $X$ est formellement lisse sur $S$, les flèches verticales sont
des équivalences de topos ;
\item en caractéristique $0$ ($S$ au-dessus de $\mathbf{Q}$), les flèches
horizontales le sont.
\end{itemize}
Les deux énoncés sont vrais, pour les raisons dites à propos de la page
2 : la lissité formelle fournit localement les rétractions, et sur
$\mathbf{Q}$ un idéal nilpotent a une et une seule structure de puissances
divisées.

Il ajoute que, dans $\mathrm{Cris}$ et $\mathrm{Cris}^{\mathrm{st}}$, les
produits fibrés et les noyaux de couples existent, et doute, avec un point
d'interrogation en marge, qu'il en soit de même dans $\mathrm{Inf}$ et
$\mathrm{Inf}^{\mathrm{st}}$.\footnote{Nous rapportons ces deux affirmations
sans les avoir vérifiées. On notera seulement que, sans puissances divisées,
le produit fibré de deux épaississements nilpotents au-dessus d'un troisième
est encore un épaississement nilpotent ; c'est avec les puissances divisées
qu'il faut passer par une enveloppe, et que la question est délicate.}

\pagerange{28}{30}
\subsection*{La suite spectrale de Čech, et la cohomologie de
$\mathcal{O}$}

Travaillons dans $\mathrm{Cris}^{\mathrm{st}}$. Soit $\widetilde{X}$ le
faisceau qui à $(U, U')$ associe l'ensemble des $S$-morphismes $U' \to X$
prolongeant l'inclusion $U \subset X$. Par définition du site stratifié,
$\widetilde{X} \to e$ est un épimorphisme sur l'objet final, et le
recouvrement qu'il définit donne une suite spectrale
\[
E_2^{p,q} = H^{p}\Bigl(\nu \mapsto H^{q}\bigl(\widetilde{X}^{\nu+1},
\mathcal{O}\bigr)\Bigr) \Longrightarrow H^{p+q}(\mathrm{Cris}^{\mathrm{st}},
\mathcal{O}),
\]
que la page appelle suite spectrale de Leray et qu'on appelle aujourd'hui
de Čech, ou de Cartan--Leray, pour un recouvrement de l'objet
final.\footnote{La page dit « l'objet $\widetilde{X}$ de
$\widetilde{\mathrm{DR}}_{\mathrm{st}}$ associé à $X$ comme l'objet final »,
la suite de la phrase étant illisible ; la description de $\widetilde{X}$
par les rétractions est la nôtre, et c'est elle qui fait de
$\widetilde{X} \to e$ un recouvrement. L'indice $\nu$ de la suite spectrale
est une lecture.}

Le produit $\widetilde{X}^{\nu+1}$ est ind-représentable : soient
$\Delta^{(\nu)}(i)$ le $i$-ième voisinage infinitésimal de la diagonale de
$X^{\nu+1}$ et $\Delta^{(\nu)}(i, j)$ son voisinage à puissances divisées
d'ordre $j$ ; alors
\[
\widetilde{X}^{\nu+1} = \varinjlim_{i, j} \Delta^{(\nu)}(i, j).
\]
Il « faudra vérifier » que, pour tout objet $(U, U')$ du site stratifié,
\[
H^{q}\bigl(\mathrm{Cris}^{\mathrm{st}}/(U, U'),\ \mathcal{O}\bigr) \simeq
H^{q}(U', \mathcal{O}_{U'}) ;
\]
c'est vrai, et c'est aujourd'hui un lemme de base de la théorie : la
cohomologie d'un faisceau sur le site cristallin localisé en un objet est la
cohomologie de Zariski de l'épaississement.\footnote{La note marginale « à
vérifier » est en regard de cette formule. La prose qui va de la formule
$\Delta^{(\nu)}(i) = \varinjlim_j \Delta^{(\nu)}(i, j)$ à
$\widetilde{X}^{\nu+1} = \varinjlim \Delta^{(\nu)}(i,j)$ est d'une écriture
très rapide et n'est que partiellement lue ; les deux formules sont sûres.}
D'où, si $X$ est affine (et séparé sur $S$, pour que la diagonale soit
fermée),
\[
H^{q}\bigl(\widetilde{X}^{\nu+1}, \mathcal{O}\bigr) = 0 \ (q > 0), \qquad
H^{0}\bigl(\widetilde{X}^{\nu+1}, \mathcal{O}\bigr) = \varprojlim_{i,j}
\Gamma\bigl(\Delta^{(\nu)}(i,j), \mathcal{O}\bigr),
\]
et par suite
\[
H^{n}(\mathrm{Cris}^{\mathrm{st}}, \mathcal{O}) \simeq H^{n}\Bigl(\nu
\mapsto H^{0}\bigl(\widetilde{X}^{\nu+1}, \mathcal{O}\bigr)\Bigr) :
\]
la cohomologie est celle d'un complexe cosimplicial explicite, formé des
anneaux des voisinages à puissances divisées des diagonales. C'est ce
qu'on appelle aujourd'hui le complexe de Čech--Alexander. Dans le cas
général, ces complexes se recollent en un complexe $\mathcal{E}^{*}$ de
faisceaux de Zariski, et $H^{*}(\mathrm{Cris}^{\mathrm{st}}, \mathcal{O})
\simeq H^{*}(X_{\mathrm{zar}}, \mathcal{E}^{*})$.\footnote{L'hypothèse
« séparé » est ajoutée. La page écrit « si $X \to$ » suivi d'un mot biffé,
lu « Aff » ; les deux lignes qui introduisent $\mathcal{E}^{*}$ sont d'une
encre presque effacée, seul l'isomorphisme final est sûr.}

\pagerange{31}{32}
\section*{L'enveloppe par générateurs et relations (pages 31 à 34)}

\subsection*{La présentation}

Retour à l'enveloppe $D_A(I)$ de la page 23, cette fois par générateurs et
relations : c'est le quotient
\[
B = A\bigl[\{\Gamma^{n}(x)\}_{x \in I,\ n \geqslant 2}\bigr] \big/ (\text{relations}),
\]
où l'on convient que $\Gamma^{1}(x) = x$, les relations étant
\[
\Gamma^{n}(x+y) = \sum_{p+q=n}\Gamma^{p}(x)\Gamma^{q}(y), \qquad
\Gamma^{n}(ax) = a^{n}\,\Gamma^{n}(x), \qquad
\Gamma^{m}(x)\Gamma^{n}(x) = \binom{m+n}{m}\Gamma^{m+n}(x),
\]
pour $x, y \in I$ et $a \in A$, et l'idéal $J$ étant engendré par les
$\Gamma^{n}(x)$.\footnote{La page écrit « $\Gamma^{n}(xy) =
x^{n}\Gamma^{n}(y)$, $x \in I$, $y \in A$ » : les rôles sont inversés, puisque
$\Gamma^{n}(y)$ n'est défini que pour $y \in I$. La relation de composition
1.4 figure d'abord dans la liste, puis est biffée : elle ne peut s'écrire
comme relation polynomiale, $\Gamma^{n}(x)$ n'étant pas un élément de $I$.}
Il reste « à vérifier que c'est une solution du problème universel » : il
faut munir $J$ de puissances divisées, et un bloc de quatre lignes, hachuré,
reprenait pour cela la construction de la page 24. C'est vrai : cette
présentation est celle de l'enveloppe (c'est, à peu de chose près,
$\Gamma_A(I)$ divisé par les $x^{[1]} - x$), et la difficulté est
exactement celle que la page désigne — la composition est héritée de
$\Gamma_A(I)$ et non imposée.

\subsection*{La simplification $p$-adique}

Supposons tous les nombres premiers $\ell \neq p$ inversibles dans $A$,
c'est-à-dire $A$ une $\mathbf{Z}_{(p)}$-algèbre. Si $n = \sum a_i p^{i}$,
$0 \leqslant a_i \leqslant p-1$, alors
\[
v_p(n!) = \sum_i a_i\, v_p\bigl((p^{i})!\bigr), \qquad v_p\bigl((p^{i})!\bigr)
= \frac{p^{i}-1}{p-1},
\]
d'où, dans toute $A$-algèbre à puissances divisées, par (1.5 bis),
\[
\gamma_n(x) = \kappa_n\, \gamma_1(x)^{a_0}\,\gamma_p(x)^{a_1}\cdots
\gamma_{p^{\nu}}(x)^{a_\nu}, \qquad \kappa_n = \frac{\prod_i
\bigl((p^{i})!\bigr)^{a_i}}{n!} \in \mathbf{Z}_{(p)}^{\times}.
\]
\footnote{La page écrit « $v_p(n) = \sum a_i v_p(p^{\nu})$ » pour
$v_p(n!) = \sum a_i\, v_p\bigl((p^{i})!\bigr)$. Elle définit $\kappa_n$ (son
$c_n$) comme la « partie première à $p$ » de $\prod((p^{i})!)^{a_i}/n!$ ; la
formule précédente montre que la valuation $p$-adique de ce rationnel est
nulle, de sorte qu'il est sa propre partie première à $p$.} Par suite $B$ est
engendré par les seuls $[x, \nu] = \Gamma^{p^{\nu}}(x)$, $\nu \geqslant 1$,
$x \in I$, les autres $\Gamma^{n}$ étant \emph{définis} par la formule
ci-dessus, et les relations deviennent : l'additivité de $\Gamma^{p^{\nu}}$
(dont le second membre ne fait intervenir que des $\Gamma^{p^{i}}$, $i <
\nu$), la relation $\Gamma^{p^{\nu}}(ax) = a^{p^{\nu}}\Gamma^{p^{\nu}}(x)$, et
les relations multiplicatives. Une note marginale demande s'il ne faut pas
prendre toutes les $\Gamma^{n}(x+y)$ ; une autre isole la relation clef,
\[
\Gamma^{p^{\nu}}(x)^{p} = q_\nu\, \Gamma^{p^{\nu+1}}(x), \qquad
q_\nu = \frac{(p^{\nu+1})!}{\bigl((p^{\nu})!\bigr)^{p}} ,
\]
la même qu'à la page 4.\footnote{La note marginale écrit
$\Gamma^{p^{\nu}}(x)^{p} = \frac{(p^{\nu+1})!}{((p^{\nu})!)^{p}\, p!}
\Gamma^{p^{\nu+1}}(x)$. Avec le $p!$ au dénominateur, c'est le coefficient de
la \emph{composée} $\Gamma^{p}(\Gamma^{p^{\nu}}(x))$ (axiome 1.4), non celui
de la \emph{puissance} $p$-ième (axiome 1.5) ; les deux diffèrent d'un
facteur $p!$. On donne la forme correcte de la puissance.}

\pagerange{32}{34}
\subsection*{Le cas d'un idéal principal}

Si $I = \lambda A$, $B$ est le quotient de $A[(\Gamma^{n}(\lambda))_{n
\geqslant 2}]$ par les relations
\[
\Gamma^{m}(\lambda)\Gamma^{n}(\lambda) - \binom{m+n}{m}\Gamma^{m+n}(\lambda),
\qquad (x^{n} - y^{n})\,\Gamma^{n}(\lambda) \quad \text{pour } x, y \in A
\text{ tels que } (x-y)\lambda = 0 .
\]
La seconde, notée $(**)$ en marge, dit que $\Gamma^{n}(\lambda x) :=
x^{n}\Gamma^{n}(\lambda)$ est bien défini ; l'additivité en résulte alors
par le binôme, et la relation $\Gamma^{n}(ax) = a^{n}\Gamma^{n}(x)$ est
triviale. Si de plus les premiers $\neq p$ sont inversibles dans $A$, il
suffit des générateurs $\Gamma^{p^{\nu}}(\lambda)$, et les relations
multiplicatives se ramènent à celles où une retenue se produit en base
$p$.\footnote{La page conclut par une formule encadrée, lourdement
surchargée, qui se lit $\Gamma^{p^{\nu}}(\lambda)^{e} =
\frac{(e\,p^{\nu+1})!}{((p^{\nu})!)^{p}}\Gamma^{e\,p^{\nu+1}}(\lambda)$
« pour $p \leqslant e < 2p-1$ ». Telle qu'elle se lit, elle ne découle pas
de 1.5, qui donne $\Gamma^{p^{\nu}}(\lambda)^{e} =
\frac{(e\,p^{\nu})!}{((p^{\nu})!)^{e}}\Gamma^{e\,p^{\nu}}(\lambda)$ ; et pour
$e = p$ on retrouve la relation clef ci-dessus. Nous ne reconstituons pas la
formule encadrée au-delà.}

Ces relations sont celles-là mêmes que la page 44 imposera, de l'autre côté,
à des éléments $\pi^{(n)}$ d'un anneau donné : ici on \emph{fabrique} un
anneau où elles valent, là on demande si l'anneau qu'on a les satisfait déjà.

\pagerange{36}{36}
\section*{Une liste d'anneaux où vaut le lemme de Poincaré (page 36)}

Soit $k$ un corps de caractéristique nulle. Le lemme de Poincaré — l'exactitude
du complexe de De Rham $\Omega^{\bullet}_{A/k}$ en degrés $> 0$, et
$H^{0} = k$ — vaut, dit la page, pour les $k$-algèbres
\[
\begin{gathered}
k[[x]][Y]/(Y^{p} - x^{q}), \qquad k[[x]][Y]/\bigl(Y^{2} + P(x)Y + Q(x)\bigr),
\qquad k[x_1, \ldots, x_n]/(x_1^{p_1}, \ldots, x_n^{p_n}), \\
k[x, y]/\mathfrak{m}^{2}, \qquad k[x,y]/(x^{3}, y^{2}, x^{2}y), \qquad
k[x,y]/(x^{3}, y^{2}, xy) .
\end{gathered}
\]
\footnote{Les exposants $p$, $q$, $p_i$ sont ici des entiers quelconques,
non le nombre premier des autres pages. Nous avons vérifié à la main les
cas $k[x]/(x^{p})$ (d'où, par Künneth, les intersections complètes
monomiales de la troisième ligne) et $k[x,y]/\mathfrak{m}^{2}$ ; les autres
sont rapportés tels que la page les donne. Les deux premiers sont des
singularités de courbes quasi-homogènes (la seconde après complétion du
carré) : c'est sous le nom de quasi-homogénéité que la question a été
reprise ensuite, par H.-J. Reiffen (1967) puis K. Saito (1971).} Et une
remarque : si $A$ est locale sur $k$, d'extension résiduelle triviale, et si
$x \in \mathfrak{m}$ vérifie $dx = 0$ dans $\Omega_{A/k}$, alors $x \in
\mathfrak{m}^{3}$.\footnote{Énoncé sans démonstration sur la page ; nous ne
l'avons pas vérifié. La référence marginale « Ann. Éc. N. Sc. 65 ou 66, J.
de Liouville » est sans autre précision, et nous ne l'identifions pas.}

\pagerange{38}{38}
\section*{Le cas d'une $\mathbf{Z}_{(p)}$-algèbre : tout est déterminé par $\gamma_p$ (pages 38 et 40)}

Supposons que les caractéristiques résiduelles de $A$ soient toutes égales
à $p > 0$, c'est-à-dire que $A$ soit une $\mathbf{Z}_{(p)}$-algèbre (et $p$
dans tous ses idéaux maximaux).\footnote{La page écrit « supposons que les
car. résiduelles de $A$ soient parfaites : $p > 0$ » ; « parfaites » est une
lecture incertaine, et ce qui suit n'utilise que l'inversibilité des premiers
$\ell \neq p$.} Soient $\gamma_m$ des puissances divisées sur un idéal de
$A$.

a) Par 1.5, $\gamma_m(x)\gamma_n(x) = \binom{m+n}{n}\gamma_{m+n}(x)$. Si $m
\equiv 0 \pmod p$ et $0 \leqslant n \leqslant p-1$, le coefficient
$\binom{m+n}{n} = (m+1)\cdots(m+n)/n!$ est premier à $p$, donc
\[
\gamma_{m+n}(x) = \binom{m+n}{n}^{-1}\gamma_m(x)\,\gamma_n(x).
\]

b) Par 1.4, $\gamma_m \circ \gamma_n = \frac{(mn)!}{(n!)^{m} m!}\,\gamma_{mn}$.
Pour $m = p$ le coefficient est $e_n = \frac{(pn)!}{(n!)^{p}\,p!}$, et
Legendre donne $v_p(e_n) = s_p(n) - 1$ ; pour $n = p$, c'est $d_m =
\frac{(pm)!}{(p!)^{m}\,m!}$, et $v_p(d_m) = 0$, de sorte que
\[
\gamma_{pm} = d_m^{-1}\; \gamma_m \circ \gamma_p .
\]
\footnote{La page note $c_n$ le coefficient $e_n$, déjà employé pour
$\kappa_n$ à la page 32 ; on renomme. Les calculs de valuation sont ceux de
la page et sont exacts.}

Ensemble, a) et b) montrent que \emph{les $\gamma_n$ sont connus dès qu'on
connaît $\gamma_p$} : écrivant $n = pm + r$ avec $0 \leqslant r < p$, on a
$\gamma_n = \binom{n}{r}^{-1}\gamma_{pm}\,\gamma_r$, avec $\gamma_r(x) =
x^{r}/r!$ et $\gamma_{pm} = d_m^{-1}\gamma_m \circ \gamma_p$, et l'on conclut
par récurrence sur $n$.

\pagerange{40}{40}
\subsection*{L'opération $\gamma_p$ seule}

L'opération $\pi = \gamma_p$ satisfait
\[
\pi(x+y) = \pi(x) + \pi(y) + \sum_{0 < i < p} \frac{x^{i}y^{p-i}}{i!\,(p-i)!},
\qquad \pi(fx) = f^{p}\,\pi(x),
\]
et aussi $p!\,\pi(x) = x^{p}$.\footnote{Cette troisième relation, cas
$n = p$ de 1.5 ter, est nécessaire et manque à la liste de la page.} Il
pose la question réciproque : une opération satisfaisant ces conditions
s'étend-elle en une structure de puissances divisées $(\gamma_n)$, avec
$\gamma_p = \pi$ ? Une tentative par récurrence sur $n$ (« Cas 1 : $N
\not\equiv 0\ (p)$, donc $N = N' + s$ \ldots ») est biffée. Il propose
ensuite la forme close
\[
\gamma_n(x) = \kappa'_n\; x^{a_0}\,\pi(x)^{a_1}\,\pi^{\circ 2}(x)^{a_2}\cdots
\pi^{\circ r}(x)^{a_r}, \qquad n = a_0 + a_1 p + \cdots + a_r p^{r},
\]
où $\pi^{\circ i}$ est l'itéré $i$-ième de $\pi$ et $\kappa'_n$ un rationnel
de valuation $p$-adique nulle. Elle est juste : $\gamma_{p^{i}}$ est, à un
facteur inversible près, $\pi^{\circ i}$ (par b), puisque les $d_m$ sont
inversibles), et l'on conclut par la formule de la page 31. La question
d'existence, elle, reste ouverte sur la page.

« Il serait intéressant de connaître ce $c_n$ mod $p$, i.e. un élément de
$\mathbf{F}_p^{\times}$. Est-ce toujours $1$ ? » La réponse est non, et elle
est simple :
\[
\kappa'_n \equiv \Bigl(\prod_i a_i!\Bigr)^{-1} \pmod p .
\]
Déjà pour $n < p$, $\kappa'_n = 1/n!$, et par exemple $\kappa'_2 = 1/2 \not\equiv
1$ dès que $p \geqslant 3$. La valeur générale découle de la congruence
classique $n!/p^{v_p(n!)} \equiv (-1)^{v_p(n!)}\prod_i a_i! \pmod p$ (qui
généralise le théorème de Wilson) appliquée aux factorielles qui composent
$\kappa'_n$ : celles des $(p^{i})!$ donnent $\pm 1$, et les signes se
compensent puisque la valuation totale est nulle. En normalisant plutôt
par $\gamma_n(x) = u_n \prod_i \gamma_{a_i}\bigl(\pi^{\circ i}(x)\bigr)$,
avec $\gamma_{a}(y) = y^{a}/a!$, on trouverait $u_n \equiv 1$.\footnote{La
réponse est de l'édition. La congruence $\kappa'_n \equiv (\prod a_i!)^{-1}$
a été vérifiée numériquement pour $p = 2, 3, 5, 7$ et $n < 200$. La page
note $c_n$ ce coefficient, troisième emploi de la lettre dans le dossier, et
écrit d'abord $\kappa'_n \in \mathbf{Z}_{p}\mathbf{Z}^{*}$, biffé, puis
« $\mathbf{Q} \cap \mathbf{Z}_p^{*}$ ».}

\pagerange{42}{42}
\section*{Applications polynomiales homogènes (page 42)}

Un feuillet à l'encre bleue clôt une note dont le début ne figure pas parmi
les pages transcrites.\footnote{La transcription situe ce début « page 40 et
avant » ; mais la page 40 transcrite traite de tout autre chose, et la page
41 est un tapuscrit étranger. Les notations $\Phi$, $M$, $N$ ne sont
introduites sur aucune page du dossier. Le lien le plus proche est la page
22, où l'algèbre $\Gamma(M)$ est caractérisée par les applications
polynomiales homogènes.} Une application polynomiale homogène $\Phi : M \to
N$ de degré $n$ définit, par
\[
\Phi\Bigl(\sum_{i=1}^{r} t_i x_i\Bigr) = \sum_{p_1 + \cdots + p_r = n}
\varphi_{p_1 \ldots p_r}(x_1, \ldots, x_r)\, t_1^{p_1}\cdots t_r^{p_r}
\]
($t_i$ indéterminées), des applications $\varphi_{p_1\ldots p_r} : M^{r} \to
N$ multihomogènes de degrés $p_1, \ldots, p_r$ — ses polarisées. Elles
satisfont les relations d'additivité obtenues en substituant des sommes,
\[
\varphi_{p_1\ldots p_r}\Bigl(\sum_{i \leqslant s_1} x_{1,i}, \ldots,
\sum_{i \leqslant s_r} x_{r,i}\Bigr) = \sum
\varphi_{(q_{k,i})}\bigl((x_{k,i})\bigr),
\]
la somme portant sur les familles $(q_{k,i})$ telles que $\sum_i q_{k,i} =
p_k$ pour chaque $k$,\footnote{La page indexe la seconde condition
$q_{1,i_2}$ où l'on attend $q_{2,i_2}$.} plus des relations de symétrie. La
phrase suivante (« On suppose que l'on a vérifié que \ldots ») s'interrompt,
et le reste du feuillet est blanc. C'est le formalisme des lois polynômes de
Roby, dont l'algèbre $\Gamma(M)$ de la page 22 est l'objet universel.

\pagerange{44}{44}
\section*{« Cristaux » : quand un idéal principal porte-t-il des puissances divisées ? (pages 44 et 46)}

Sous ce titre, de sa main, une question précise. Soit $J = \pi A$ un idéal
principal. Formellement, des puissances divisées sur $J$ seraient
\[
\gamma_n(\pi y) = \frac{(\pi y)^{n}}{n!} = \frac{\pi^{n}}{n!}\, y^{n} =
\pi^{(n)} y^{n} ;
\]
on est donc amené à chercher des éléments $\pi^{(n)} \in \pi A$, $n
\geqslant 1$, tels que
\[
(1) \qquad
\begin{cases}
\pi^{(1)} = \pi, \\
\pi^{(m)}\pi^{(n)} = \dbinom{m+n}{m}\pi^{(m+n)}, \\
\pi^{(a)}\Bigl(\dfrac{\pi^{(b)}}{\pi}\Bigr)^{a} = \dfrac{(ab)!}{a!\,(b!)^{a}}\,
\pi^{(ab)} .
\end{cases}
\]
\footnote{La page écrit la troisième condition avec les lettres $p$, $q$,
qui sont ici des entiers quelconques et non le nombre premier ; on les
remplace par $a$, $b$. Un premier état des conditions, cerné et biffé,
précède.} Comme $\pi^{(b)} \in \pi A$, le quotient $\pi^{(b)}/\pi$ n'est
défini qu'à l'addition près d'un $z$ tel que $\pi z = 0$ ; mais comme
$\pi^{(a)} \in \pi A$ aussi, le produit $\pi^{(a)}(\pi^{(b)}/\pi)^{a}$ est
bien déterminé. \emph{Alors $\gamma_n(\pi y) = \pi^{(n)}y^{n}$ définit une
structure de puissances divisées sur $J$}, indépendante du choix de $y$ pour
la même raison. C'est exact : la deuxième condition donne 1.2 et 1.5, la
troisième 1.4, et 1.3 est évidente.

\emph{Exemple 1.} Si $A$ est sans torsion sur $\mathbf{Z}$, les relations
$\pi^{n} = n!\,\pi^{(n)}$ déterminent les $\pi^{(n)}$, et les conditions (1)
sont alors automatiques. Leur existence équivaut à
\[
\pi^{n} \in n!\,\pi A \qquad \text{pour tout } n \geqslant 2 .
\]
\footnote{La page écrit « $\pi^{n} \in n!\,A$ ». Ce n'est pas suffisant : il
faut $\pi^{(n)} = \pi^{n}/n! \in \pi A$ et non seulement dans $A$. Dans
l'algèbre à puissances divisées $A = \mathbf{Z}_{(p)}\langle T \rangle$,
sans torsion, avec $\pi = T$, on a $\pi^{n} \in n!\,A$ pour tout $n$, mais
$T^{p}/p! = T^{[p]}$ n'est pas dans $TA$ : l'idéal \emph{principal} $TA$ ne
porte pas de puissances divisées (c'est l'idéal qu'elles engendrent qui en
porte). Dans le cas d'un anneau de valuation discrète, qui est celui de
l'exemple 3, les deux conditions coïncident.}

\emph{Exemple 2.} Pour $A_0 = \mathbf{Z}_{(p)}$ et $\pi = p$, les $p^{n}/n!$
sont des entiers $p$-adiques, et même des éléments de $p\mathbf{Z}_{(p)}$ ;
les conditions (1) sont satisfaites. Par spécialisation, toute
$\mathbf{Z}_{(p)}$-algèbre $A$ — tout anneau où les premiers $\ell \neq p$
sont inversibles — reçoit les $\pi^{(n)} = p^{n}/n!$, qui satisfont (1) :
d'où une structure de puissances divisées canonique sur $pA$.\footnote{La
phrase commence en bas de la page 44 et s'achève en haut de la page 46 ; la
page 45, son verso, porte une page du tapuscrit d'EGA donné plus bas.
« inversible » et « caractéristiques » sont des lectures incertaines.} C'est
la structure que l'on suppose aujourd'hui, sous le nom de compatibilité, sur
la base $\mathbf{Z}_p$ de la cohomologie cristalline ; et ce sont les
puissances divisées $\lambda_i = \gamma_{p^{i}}(p)$ de la page 4, celles dont
l'anneau universel $W^{(\nu)}$ mesurait l'écart.

\pagerange{46}{46}
\emph{Exemple 3.} Soit $A$ intègre, de caractéristiques résiduelles $p$ et
$0$, et $\pi$ tel que $A/\pi A$ soit de caractéristique $p$. Quand
existe-t-il des $\pi^{(n)}$ ? Supposons $A$ de valuation discrète, $v$ sa
valuation normalisée, $e = v(p)$ son indice de ramification absolu, de sorte
que $v(N) = e\,v_p(N)$ pour un entier $N$. Il faut $\pi^{n}/n! \in A$,
c'est-à-dire
\[
v(\pi) \geqslant e\,\frac{v_p(n!)}{n} \quad \text{pour tout } n, \qquad
\text{soit} \qquad v(\pi) \geqslant e \sup_n \frac{v_p(n!)}{n} =
\frac{e}{p-1},
\]
la borne supérieure valant $1/(p-1)$ par Legendre (nom écrit en marge) —
elle n'est pas atteinte, mais c'est la condition pour tout $n$ qui compte.
Si $\pi$ est une uniformisante, $v(\pi) = 1$, d'où la condition
\[
\boxed{\ e \leqslant p-1\ }
\]
C'est, sous sa forme définitive, le critère de la page 12 : l'idéal maximal
d'un anneau de valuation discrète de caractéristique mixte porte des
puissances divisées si et seulement si $e \leqslant p - 1$ (et elles sont
alors uniques). Il gouvernera, dans toute la suite de l'histoire, la
frontière entre ce que la théorie cristalline sait faire sans peine et ce
qu'elle ne sait faire qu'avec des efforts — les théorèmes de comparaison de
Fontaine--Messing, par exemple, supposent $e = 1$ et $p$ assez grand.\footnote{La
page dit « $\lambda(p) = \sup_n v_p(n!)/n = 1/(p-1)$ est une constante ne
dépendant que de $p$, et qui est $\leqslant 1$ ». La mention de
Fontaine--Messing est de l'édition. Que la condition soit aussi suffisante
résulte du calcul de la page 12, où $v(\pi^{n}/n!) \geqslant v(\pi)$ dès que
$v_p(\pi) \geqslant 1/(p-1)$, de sorte que $\pi^{n}/n! \in \pi A$.}

\pagerange{45}{47}
\section*{Un tapuscrit d'EGA IV, § 18.13, annoté (pages 45 et 47)}

Au dos des deux feuillets « Cristaux », deux pages d'un tapuscrit d'EGA IV,
numérotées 22 et 23 du § 18.13, qui se lisent dans cet ordre (page 47, puis
page 45).\footnote{La page 22 du tapuscrit finit sur « co- » et la page 23
commence par « ïncide ». Nous n'avons pas comparé ce tapuscrit au texte
imprimé d'EGA IV.} Leur objet est le \emph{module des différentielles
complété} d'une algèbre topologique.

Soit $A$ un anneau topologique et $B$ une $A$-algèbre topologique qui est un
anneau semi-local noethérien muni de sa topologie $\mathfrak{r}$-préadique
($\mathfrak{r}$ son radical). Pour tout $B$-module de type fini $L$, toute
$A$-dérivation $D : B \to L$ est continue, puisque $D(\mathfrak{r}^{j+1})
\subset \mathfrak{r}^{j}L$ ; elle se factorise donc par $\Omega^{1}_{B/A}$.
On définit $\overline{\Omega}^{1}_{B/A}$ comme le quotient de
$\Omega^{1}_{B/A}$ par l'intersection des noyaux de ses applications vers
des $B$-modules de type fini, et $\bar d_{B/A}$ la dérivation composée ;
alors
\[
(18.13.4.1) \qquad \mathrm{Hom}_B\bigl(\overline{\Omega}^{1}_{B/A}, L\bigr)
\xrightarrow{\ \sim\ } \mathrm{D\acute{e}r}_A(B, L) =
\mathrm{D\acute{e}r.cont}_A(B, L)
\]
pour tout $B$-module de type fini $L$. Si $\overline{\Omega}^{1}_{B/A}$ est de
type fini, il représente donc le foncteur $L \rightsquigarrow
\mathrm{D\acute{e}r}_A(B, L)$ sur les modules de type fini, et si
$\Omega^{1}_{B/A}$ lui-même est de type fini (par exemple si $B$ est de type
fini sur $A$), les deux coïncident. La proposition (18.13.5) traite des
suites exactes de ces modules pour $A \to B \to C$.

Sa main ajoute le renvoi $(0_{\mathrm{III}}, 8.1.11)$, l'égalité
« $= \mathrm{D\acute{e}r.cont}_A(B, L)$ », la note marginale
« On a donc $\mathrm{Hom}_B(\overline{\Omega}^{1}_{B/A}, L) =
\mathrm{Hom}_B(\Omega^{1}_{B/A}, L)$ pour tout $B$-module de type fini $L$ »,
biffe dans (18.13.5) les mots « déterminés de façon unique », et pose en
marge, à côté de « semi-local noethérien », la question d'un rédacteur :
« à quoi sert l'hypothèse ? Prendre topologie plus générale ? ». Ces pages
ne touchent pas aux cristaux ; elles disent seulement de quoi était faite sa
table de travail.

\pagerange{49}{50}
\section*{La lettre à Deligne du 10 décembre 1965 (pages 49 à 52)}

Une chemise de sa main, « lettre à Deligne / sur coh. $p$-adique », couvre le
double dactylographié d'une lettre datée à la main.\footnote{Le double n'a
pas les flèches, restituées d'après le texte, ni les crochets de décalage
$[2d]$.}

\subsection*{Le théorème de dualité}

La lettre s'ouvre sur une simplification de la démonstration du théorème
de dualité pour la cohomologie étale (celui de SGA 4, exposé XVIII), qui
évite le théorème de pureté relative. On est ramené au cas où
$f : X \to Y$ est lisse de dimension relative $1$, les coefficients étant
constants ; le passage à la limite de SGA 4, exposé VI, § 6, permet de
supposer la base noethérienne et de dimension finie, et l'on raisonne par
récurrence sur la dimension : localisant sur $Y$, on suppose $Y$ strictement
local de point fermé $y$, et l'on traite séparément $X \times_Y (Y - y)$,
par hypothèse de récurrence, et $X \times_Y y$.\footnote{La lettre fait la
récurrence sur « $n = \dim X$ » et invoque « $\dim(Y - y) < n$ » ; c'est
la dimension qui baisse quand on retire le point fermé de la base qui fait
marcher l'argument.}

Il ajoute que le théorème vaut, avec le même énoncé et la même preuve, pour
un morphisme lisse compactifiable $f : (X, B) \to (Y, A)$ d'espaces annelés,
où $A$ est un faisceau d'anneaux sur $Y$ annulé par un entier $n$ premier aux
caractéristiques résiduelles, et $B$ un faisceau d'anneaux sur $X$ muni de
$f^{-1}(A) \to B$, en posant, pour $d$ la dimension relative,
\[
f^{!}(K) = R\,\underline{\mathrm{Hom}}_{f^{-1}(A)}\bigl(B,\ f^{-1}(K)
\otimes T_{X/Y}[2d]\bigr),
\]
où $T_{X/Y}$ est le faisceau de Tate (les racines $n$-ièmes de l'unité,
tordues $d$ fois). La trace et la dualité se décomposent alors entre le cas
où $f^{-1}(A) \to B$ est un isomorphisme, qui se traite comme le cas
$A = (\mathbf{Z}/n\mathbf{Z})_Y$, et le cas $f = \mathrm{id}$, trivial. « Je
pense qu'il vaut le coup d'inclure cette forme générale du théorème de
dualité, soit de prime abord, soit à la fin dans un
numéro-ratrappage. »\footnote{« numéro-ratrappage » est son orthographe,
vérifiée sur le fac-similé.} Il doute qu'elle se déduise du cas constant
comme simple corollaire, et la même remarque vaut pour la dualité globale.

\pagerange{50}{51}
\subsection*{Une comparaison $p$-adique}

Vient le passage qui donne son titre à la chemise. Sur $\mathbf{C}$, la
cohomologie de De Rham et la cohomologie $\ell$-adique sont reliées, de
façon transcendante, par la cohomologie entière. Il devrait y avoir un
analogue sur un corps $K$ algébriquement clos, complet pour une valeur
absolue non archimédienne, de caractéristique $0$ et de corps résiduel de
caractéristique $p$ : une application $\mathbf{Z}_p$-linéaire
$H^{*}(X, \mathbf{Z}_p) \to H^{*}_{\mathrm{DR}}(X)$ qui induirait un
isomorphisme $H^{*}(X, \mathbf{Z}_p) \otimes_{\mathbf{Z}_p} K \simeq
H^{*}_{\mathrm{DR}}(X)$.

Pour la construire, il propose de passer par l'espace rigide-analytique
$X^{\mathrm{an}}$ de Tate — « un site qui ne peut être défini par une
topologie au sens ordinaire » —, d'espérer que sa cohomologie à coefficients
de torsion coïncide avec la cohomologie étale de $X$ et, « par un excès
d'optimisme », que sa cohomologie à coefficients constants $\mathbf{Z}_p$
soit la cohomologie $p$-adique. L'inclusion de $\mathbf{Z}_p$ dans les
fonctions enverrait alors $H(X^{\mathrm{an}}, \mathbf{Z}_p)$ dans la
cohomologie de De Rham de $X^{\mathrm{an}}$, qui est celle de $X$ par un
GAGA rigide. Il faudrait pour cela un « lemme de Poincaré »
rigide-analytique.

Et aussitôt il doute : il n'est pas sûr que la cohomologie rigide à
coefficients constants donne les bons nombres de Betti, et se rappelle que
Tate lui a dit que c'était déjà faux pour le $H^{1}$ d'une courbe
elliptique.

Les deux doutes étaient fondés, et le premier va plus loin qu'il ne le
dit.\footnote{Ce paragraphe est de l'édition.} La cohomologie à
coefficients constants de l'analytifiée ne voit que la topologie de
l'espace analytique — dans le langage de Berkovich, venu plus tard, elle
est de rang $0$ en degré $1$ pour une courbe elliptique à bonne réduction
et de rang $1$ pour une courbe de Tate, jamais $2$. Et, plus profondément,
il n'existe pas d'isomorphisme naturel $H^{*}_{\text{ét}}(X,
\mathbf{Z}_p) \otimes K \simeq H^{*}_{\mathrm{DR}}(X)$ : il faut, pour
comparer les deux, un anneau de « périodes » $p$-adiques, $B_{\mathrm{dR}}$,
construit par Fontaine au début des années 1980, et le théorème de
comparaison — démontré par Faltings, Tsuji et d'autres, puis étendu par
Scholze aux espaces rigides — dit $H^{*}_{\text{ét}} \otimes
B_{\mathrm{dR}} \simeq H^{*}_{\mathrm{DR}} \otimes B_{\mathrm{dR}}$. C'est le
« foncteur mystérieux » qu'il appellera de ses vœux en 1970. Le lemme de
Poincaré rigide, lui, au niveau des faisceaux, ne pose pas de difficulté :
une forme fermée sur un polydisque s'intègre sur tout polydisque de rayon
strictement plus petit, et ces polydisques recouvrent.

\pagerange{51}{52}
\subsection*{En caractéristique $p$}

En caractéristique $p > 0$, au contraire, il y a une flèche évidente de
$H^{*}(X, \mathbf{F}_p)$ dans la cohomologie de De Rham — et c'est le point
de départ de la suite. On calcule la cohomologie de De Rham sur le site
étale (le résultat est le même que sur le site de Zariski, les composantes
du complexe étant quasi-cohérentes), au moyen de la suite spectrale
conjuguée
\[
E_2^{p,q} = H^{p}\bigl(X, \mathcal{H}^{q}(\Omega^{\bullet})\bigr)
\Longrightarrow H^{p+q}_{\mathrm{DR}}(X).
\]
La flèche se factorise par $H^{*}(X, \mathbf{F}_p) \to H^{*}(X,
\mathcal{O}_X)$, puis $H^{*}(X, \mathcal{O}_X) \to H^{*}_{\mathrm{DR}}(X)$,
cette dernière venant de l'isomorphisme $\mathcal{O}_X \simeq
\mathcal{H}^{0}(\Omega^{\bullet})$, $f \mapsto f^{p}$.\footnote{Hypothèse
implicite : $X$ lisse sur un corps parfait. C'est alors le degré $0$ de
l'isomorphisme de Cartier ; sans lissité, $\mathcal{H}^{0}(\Omega^{\bullet})$
n'est pas en général $\mathcal{O}_X^{p}$.} Le composé
\[
H^{n}(X, \mathcal{O}_X) \longrightarrow H^{n}_{\mathrm{DR}}(X)
\longrightarrow H^{n}(X, \mathcal{O}_X),
\]
la seconde flèche venant de l'autre suite spectrale (celle de Hodge vers De
Rham), « ne peut guère être autre chose que l'homomorphisme de
Frobenius » : c'est bien lui, puisque la première flèche est induite par
$f \mapsto f^{p}$ et que la seconde est l'arête $H^{n}_{\mathrm{DR}} \to
H^{n}(X, \mathcal{O}_X)$.

« Il faut dire que tout ça est bien éculé, et qu'on ne pourra dire des choses
vraiment intéressantes et nouvelles qu'en faisant appel à la “vraie”
cohomologie $p$-adique. » Un an plus tard, à l'IHES, il exposera les
cristaux ; les feuillets qui précèdent cette lettre dans le dossier en sont
l'établi.

\end{document}
